%%%%%%%%%%%%%%%%%%%%%%%%%%%%%%%%%%%%%%%%%%%%%%%%%%%%%%%%%%%%%%%%%%%%%%%%%%%%%%
%%%%%%%%%%%%%%%%%%%%%%%%%%%%%%%%%%%%%%%%%%%%%%%%%%%%%%%%%%%%%%%%%%%%%%%%%%%%%%
%%%%%%%%%%%%%%%%%%%%%%%%%%%%%%%%%%%%%%%%%%%%%%%%%%%%%%%%%%%%%%%%%%%%%%%%%%%%%%

\chapter{Fungsi sebagai Limit} \label{approx:chapter}

%%%%%%%%%%%%%%%%%%%%%%%%%%%%%%%%%%%%%%%%%%%%%%%%%%%%%%%%%%%%%%%%%%%%%%%%%%%%%%

\section{Bilangan kompleks}
\label{sec:complexnums}

%mbxINTROSUBSECTION

\sectionnotes{setengah kuliah}

\subsection{Bidang kompleks}

Dalam bab ini, kita membahas aproksimasi fungsi, atau dengan kata lain,
fungsi sebagai limit barisan dan deret.
Kita akan memperluas beberapa hasil yang telah kita lihat ke suatu
latar yang sedikit lebih umum, dan akan melihat beberapa hasil yang sama sekali baru.
Secara khusus, kita membahas fungsi bernilai kompleks.
Sebelumnya kita telah memberikan bilangan kompleks sebagai contoh, tetapi
mari kita mulai dari awal dan mendefinisikan lapangan bilangan kompleks secara tepat.

Bilangan kompleks hanyalah pasangan $(x,y) \in \R^2$ yang padanya kita
mendefinisikan perkalian (lihat di bawah).
Kita menyebut himpunan ini \emph{bilangan kompleks}\index{bilangan kompleks}
dan menotasikannya dengan $\C$.
Kita mengidentifikasi $x \in \R$ dengan $(x,0) \in \C$.
Sumbu-$x$ kemudian disebut \emph{\myindex{sumbu real}} dan sumbu-$y$ disebut
\emph{\myindex{sumbu imajiner}}.
Karena $\C$ hanyalah bidang, kita juga menyebut himpunan $\C$ sebagai
\emph{\myindex{bidang kompleks}}.

Definisikan:
\begin{equation*}
(x,y) + (s,t) \coloneqq (x+s,y+t) , \qquad
(x,y) (s,t) \coloneqq (xs-yt,xt+ys) .
\end{equation*}
Dengan identifikasi di atas, kita memiliki $0 = (0,0)$ dan $1 = (1,0)$.
Kedua operasi ini menjadikan bidang tersebut sebuah lapangan (latihan).
Kita menulis bilangan kompleks $(x,y)$ sebagai $x+iy$, dengan
mendefinisikan\footnote{Perhatikan bahwa para insinyur menggunakan $j$ sebagai pengganti $i$.}
\glsadd{not:imaginary}
\begin{equation*}
i \coloneqq (0,1) .
\end{equation*}
Perhatikan bahwa $i^2 = (0,1)(0,1) = (0-1,0+0) = -1$.
Artinya, $i$ adalah solusi persamaan polinomial
\begin{equation*}
z^2+1=0 .
\end{equation*}
Mulai sekarang, kita tidak akan menggunakan notasi $(x,y)$ dan hanya menggunakan $x+iy$.
Lihat \figureref{fig:complexplane}.
\begin{myfigureht}
\myincludegraphics{complexplane}{%
Diagram sebuah bilangan pada bidang kompleks. Bilangan 1 ditandai pada
sumbu horizontal di sebelah kanan titik asal,
dan bilangan i ditandai pada sumbu vertikal di atas titik asal.
Bilangan x ditandai pada sumbu horizontal dan bilangan i y ditandai
pada sumbu vertikal. Kombinasi x ditambah i y ditandai pada titik
(x,y).}
\caption{Titik-titik $1$, $i$, $x$, $iy$, dan $x+iy$ pada bidang
kompleks.\label{fig:complexplane}}
\end{myfigureht}

Kita umumnya menggunakan $x,y,r,s,t$ untuk nilai real dan $z,w,\xi,\zeta$
untuk nilai kompleks, meskipun ini bukan aturan yang kaku. Khususnya,
$z$ sering digunakan sebagai variabel real ketiga dalam $\R^3$.

\begin{defn}
Misalkan $z= x+iy$.
Kita menyebut $x$ 
sebagai \emph{\myindex{bagian real}} dari $z$, dan 
menyebut $y$
sebagai \emph{\myindex{bagian imajiner}} dari $z$. Kita menulis
\glsadd{not:realpart}\glsadd{not:imagpart}
\begin{equation*}
\Re\, z \coloneqq x , \qquad
\Im\, z \coloneqq y .
\end{equation*}
Definisikan
\glsadd{not:conj}
\emph{\myindex{konjugat kompleks}} sebagai
\begin{equation*}
\bar{z} \coloneqq x-iy ,
\end{equation*}
dan definisikan \emph{\myindex{modulus}} sebagai
\glsadd{not:modulus}
\begin{equation*}
\sabs{z} \coloneqq \sqrt{x^2+y^2} .
\end{equation*}
\end{defn}

Modulus merupakan analog kompleks dari nilai mutlak dan
memiliki sifat-sifat serupa.
Sebagai contoh,
$\sabs{zw} = \sabs{z} \, \sabs{w}$ (latihan).
Konjugat kompleks merupakan pencerminan bidang terhadap sumbu real.
Bilangan real tepatnya adalah bilangan-bilangan yang bagian imajinernya
$y=0$. Secara khusus, bilangan-bilangan tersebut tepatnya adalah
bilangan yang memenuhi persamaan
\begin{equation*}
z = \bar{z} .
\end{equation*}

Karena $\C$ sebenarnya adalah $\R^2$, kita gunakan pada $\C$ metrik
Euklides standar pada $\R^2$.
Secara khusus,
\begin{equation*}
\sabs{z} = d(z,0) , \qquad 
\text{dan juga} \qquad 
\sabs{z-w} = d(z,w) .
\end{equation*}
Jadi, topologi pada $\C$ sama persis dengan topologi standar pada $\R^2$
dengan metrik Euklides,
dan $\sabs{z}$ sama dengan norma Euklides pada $\R^2$.
Yang penting, karena $\R^2$ merupakan ruang metrik lengkap, maka
$\C$ juga demikian.
Karena $\sabs{z}$ merupakan norma Euklides pada $\R^2$, kita memiliki
\emph{ketaksamaan segitiga}\index{ketaksamaan segitiga!bilangan kompleks}
dalam kedua bentuk:
\begin{equation*}
\sabs{z+w} \leq \sabs{z}+\sabs{w} \qquad \text{dan} \qquad
\big\lvert \sabs{z}-\sabs{w} \big\rvert \leq \sabs{z-w} .
\end{equation*}

Konjugat kompleks dan modulus ternyata berkaitan lebih erat:
\begin{equation*}
\sabs{z}^2 =
x^2+y^2 =
(x+iy)(x-iy) =
z \bar{z} .
\end{equation*}

\begin{remark}
Tidak ada urutan alami pada bilangan kompleks.
Secara khusus,
tidak ada urutan yang menjadikan bilangan kompleks sebagai lapangan terurut.
Pengurutan adalah salah satu hal yang hilang ketika kita beralih dari bilangan real
ke bilangan kompleks.
\end{remark}

\subsection{Bilangan kompleks dan limit}

Operasi aljabar pada bilangan kompleks bersifat kontinu karena konvergensi di
$\R^2$ sama dengan konvergensi pada setiap komponen, dan kita telah mengetahui
bahwa operasi aljabar real bersifat kontinu. Sebagai contoh,
tuliskan $z_n = x_n + i\,y_n$ dan
$w_n = s_n + i\,t_n$, dan misalkan
$\lim_{n\to\infty} z_n = z = x+i\,y$ serta $\lim_{n\to\infty} w_n = w = s+i\,t$.
Mari kita tunjukkan
\begin{equation*}
\lim_{n\to\infty} z_n w_n = zw .
\end{equation*}
Pertama,
\begin{equation*}
z_n w_n = (x_ns_n-y_nt_n) + i(x_nt_n+y_ns_n) .
\end{equation*}
Topologi pada $\C$ sama dengan topologi pada $\R^2$, sehingga
$x_n \to x$, $y_n \to y$, $s_n \to s$, dan $t_n \to t$.
Dengan demikian,
\begin{equation*}
\lim_{n\to\infty} (x_ns_n-y_nt_n) = xs-yt \qquad \text{dan} \qquad
\lim_{n\to\infty} (x_nt_n+y_ns_n) = xt+ys .
\end{equation*}
Karena
$(xs-yt)+i(xt+ys) = zw$,
\begin{equation*}
\lim_{n\to\infty} z_n w_n = zw .
\end{equation*}

Demikian pula, modulus dan konjugat kompleks merupakan fungsi kontinu. Kita
serahkan sisa pembuktian proposisi berikut sebagai latihan.

\begin{prop} \label{prop:continuityofcomplex}
Misalkan $\{ z_n \}_{n=1}^\infty$, $\{ w_n \}_{n=1}^\infty$ merupakan barisan bilangan kompleks yang masing-masing konvergen
ke $z$ dan $w$. Maka
\begin{enumerate}[(i)]
\item
$\displaystyle \lim_{n\to \infty} (z_n + w_n) = z + w$.
\item
$\displaystyle \lim_{n\to \infty} z_n w_n = z w$.
\item
Andaikan $w_n \neq 0$ untuk semua $n$ dan $w \neq 0$,
$\displaystyle \lim_{n\to \infty} \frac{z_n}{w_n} = \frac{z}{w}$.
\item
$\displaystyle \lim_{n\to \infty} \sabs{z_n} = \sabs{z}$.
\item
$\displaystyle \lim_{n\to \infty} \bar{z}_n = \bar{z}$.
\end{enumerate}
\end{prop}

Seperti telah kita lihat di atas, konvergensi di $\C$ sama dengan konvergensi di
$\R^2$. Secara khusus, suatu barisan di $\C$ konvergen jika dan hanya jika
bagian real dan imajinernya konvergen. Karena itu, silakan terapkan
semua yang telah Anda pelajari tentang konvergensi di $\R^2$, serta
terapkan hasil-hasil tentang bilangan real pada bagian real dan imajinernya.

Kita juga memerlukan konvergensi deret kompleks. Misalkan $\{ z_n \}_{n=1}^\infty$ merupakan
barisan bilangan kompleks. Deret
\begin{equation*}
\sum_{n=1}^\infty z_n
\end{equation*}
\emph{konvergen}\index{konvergen!deret kompleks} jika barisan jumlah parsialnya konvergen, yaitu jika
\begin{equation*}
\lim_{k\to\infty} \sum_{n=1}^k z_n \qquad \text{ada.}
\avoidbreak
\end{equation*}
Suatu deret \emph{konvergen absolut}\index{konvergen absolut!deret kompleks}
jika $\sum_{n=1}^\infty \sabs{z_n}$ konvergen.

Kita mengatakan suatu deret
bersifat \emph{Cauchy}\index{Cauchy!deret kompleks}
jika barisan jumlah parsialnya bersifat Cauchy. Dua
proposisi berikut pada dasarnya memiliki pembuktian yang sama
dengan padanannya untuk deret real, dan
kita serahkan sebagai latihan.

\begin{prop} \label{prop:cachysercomplex}
Deret kompleks $\sum_{n=1}^\infty z_n$ bersifat Cauchy jika untuk setiap $\epsilon > 0$
terdapat $M \in \N$ sedemikian sehingga untuk setiap $n \geq M$
dan setiap $k > n$, berlaku
\begin{equation*}
\abs{ \sum_{j={n+1}}^k z_j }
< \epsilon .
\end{equation*}
\end{prop}

\begin{prop} \label{prop:absconvmeansconv}
Jika suatu deret kompleks $\sum_{n=1}^\infty z_n$ konvergen absolut, maka deret tersebut konvergen.
\end{prop}

Deret $\sum_{n=1}^\infty \sabs{z_n}$ merupakan deret real. Semua
uji konvergensi (uji rasio, uji akar, dan sebagainya)\ yang membahas
konvergensi absolut bekerja pada bilangan $\sabs{z_n}$; artinya, uji-uji tersebut
sebenarnya membahas konvergensi deret bilangan real nonnegatif.
Anda dapat langsung menerapkan uji-uji ini
tanpa perlu membuktikan ulang apa pun untuk deret kompleks.

\subsection{Fungsi bernilai kompleks}

Ketika membahas fungsi bernilai kompleks
$f \colon X \to \C$, kita sering menulis
$f = u+i\,v$ untuk fungsi bernilai real $u \colon X \to \R$ dan $v \colon X \to
\R$.

Misalkan kita ingin mengintegralkan
$f \colon [a,b] \to \C$. Kita menulis
$f = u+i\,v$ untuk $u$ dan~$v$ yang bernilai real.
Kita mengatakan bahwa $f$ \emph{terintegralkan secara Riemann}\index{terintegralkan secara Riemann!
fungsi bernilai kompleks}
jika $u$ dan $v$ terintegralkan secara Riemann,
dan dalam hal ini kita mendefinisikan
\begin{equation*}
\int_a^b f \coloneqq \int_a^b u + i \int_a^b v .
\end{equation*}
Kita menggunakan definisi yang sama untuk setiap jenis integral lainnya (tak wajar,
multivariabel, dan sebagainya).

Demikian pula, ketika melakukan diferensiasi, tuliskan $f \colon [a,b] \to \C$ sebagai
$f = u+i\,v$. Dengan memandang $\C$ sebagai $\R^2$, kita mengatakan bahwa $f$ diferensiabel
jika $u$ dan $v$ diferensiabel. Untuk fungsi yang bernilai di $\R^2$, turunannya
direpresentasikan oleh suatu vektor di $\R^2$. Vektor di $\R^2$ tersebut sekarang merupakan bilangan
kompleks. Dengan kata lain,
kita menulis
\emph{turunan}\index{turunan!fungsi bernilai kompleks}
sebagai
\glsadd{not:mvder}
\begin{equation*}
f'(t) \coloneqq u'(t) + i \, v'(t) .
\end{equation*}
Operator linear yang merepresentasikan turunan adalah perkalian dengan
bilangan kompleks $f'(t)$, sehingga tidak ada yang hilang dalam identifikasi ini.


\subsection{Latihan}

\begin{exercise}
Periksa bahwa $\C$ merupakan lapangan.
\end{exercise}

\begin{exercise}
Buktikan bahwa untuk $z,w \in \C$ berlaku
$\sabs{zw} = \sabs{z} \, \sabs{w}$.
\end{exercise}

\begin{exercise}
Selesaikan pembuktian \propref{prop:continuityofcomplex}.
\end{exercise}

\begin{exercise}
Buktikan \propref{prop:cachysercomplex}.
\end{exercise}

\begin{exercise}
Buktikan \propref{prop:absconvmeansconv}.
\end{exercise}

\begin{samepage}
\begin{exercise}
Untuk $x +iy$, definisikan matriks
$\left[ \begin{smallmatrix} x & -y \\ y & x \end{smallmatrix} \right]$.
Buktikan:
\begin{enumerate}[a)]
\item
Aksi matriks ini pada vektor $(s,t)$ sama dengan mengalikan
$(x+iy)(s+it)$.
\item
Mengalikan dua matriks semacam ini sama dengan mengalikan bilangan kompleks
yang mendasarinya, lalu mencari matriks yang bersesuaian dengan hasil kali tersebut.
Dengan kata lain, lapangan $\C$ dapat diidentifikasi dengan suatu subhimpunan matriks 2-kali-2.
\item
Matriks
$\left[ \begin{smallmatrix} x & -y \\ y & x \end{smallmatrix} \right]$
memiliki nilai eigen $x+iy$ dan $x-iy$. Ingat bahwa $\lambda$ adalah
nilai eigen matriks $A$ jika $A-\lambda I$ (matriks kompleks dalam kasus kita)
tidak invertibel, yaitu,
jika baris-barisnya bergantung secara linear: satu baris merupakan kelipatan
(kompleks) dari baris lainnya.
\end{enumerate}
\end{exercise}
\end{samepage}

\begin{exercise}
Buktikan teorema Bolzano--Weierstrass untuk barisan kompleks.
Misalkan $\{ z_n \}_{n=1}^\infty$ merupakan barisan bilangan kompleks yang terbatas.
Artinya, terdapat $M$ sedemikian sehingga $\sabs{z_n} \leq M$ untuk semua $n$.
Buktikan bahwa terdapat subbarisan $\{ z_{n_k} \}_{k=1}^\infty$
yang konvergen ke suatu $z \in \C$.
\end{exercise}

\begin{exercise}
\leavevmode
\begin{enumerate}[a)]
\item
Buktikan bahwa tidak ada teorema nilai rata-rata sederhana untuk fungsi bernilai kompleks:
Temukan fungsi diferensiabel $f \colon [0,1] \to \C$ sedemikian sehingga
$f(0) = f(1) = 0$, tetapi $f'(t) \neq 0$ untuk semua $t \in [0,1]$.
\item
Namun, ada bentuk teorema nilai rata-rata yang lebih lemah sebagaimana untuk fungsi bernilai vektor.
Buktikan: Jika $f \colon [a,b] \to \C$ kontinu dan diferensiabel pada
$(a,b)$, dan untuk suatu $M$ berlaku $\babs{f'(x)} \leq M$ untuk semua $x \in (a,b)$, maka
$\babs{f(b)-f(a)} \leq M \sabs{b-a}$.
\end{enumerate}
\end{exercise}

\begin{exercise}
Buktikan bahwa tidak ada teorema nilai rata-rata sederhana untuk integral
fungsi bernilai kompleks:
Temukan fungsi kontinu $f \colon [0,1] \to \C$ sedemikian sehingga
$\int_0^1 f = 0$ tetapi $f(t) \neq 0$ untuk semua $t \in [0,1]$.
\end{exercise}


%%%%%%%%%%%%%%%%%%%%%%%%%%%%%%%%%%%%%%%%%%%%%%%%%%%%%%%%%%%%%%%%%%%%%%%%%%%%%%

\sectionnewpage
\section{Pertukaran limit}
\label{sec:swaplim}

%mbxINTROSUBSECTION

\sectionnotes{2 kuliah}

\subsection{Kekontinuan}

Mari kita kembali ke pertukaran limit dan memperluas pembahasan
\volIref{\chapterref*{vI-fs:chapter} dari jilid I}{\chapterref{fs:chapter}}.
Misalkan $\{ f_n \}_{n=1}^\infty$ suatu barisan
fungsi $f_n \colon X \to Y$ untuk suatu himpunan $X$ dan ruang metrik $Y$.
Misalkan $f \colon X \to Y$ suatu
fungsi dan andaikan untuk setiap $x \in X$ berlaku
\begin{equation*}
f(x) = \lim_{n\to \infty} f_n(x) .
\end{equation*}
Kita mengatakan barisan $\{ f_n \}_{n=1}^\infty$
\emph{\myindex{konvergen titik demi titik}}\index{konvergensi titik demi titik} ke $f$.

Untuk $Y=\C$, suatu deret fungsi
\emph{konvergen titik demi titik}\index{konvergen titik demi titik!deret kompleks}\index{konvergensi titik demi titik!deret kompleks} ke $f$ jika
untuk setiap $x \in X$, berlaku
\begin{equation*}
f(x) = \lim_{n\to \infty} \sum_{k=1}^n f_k(x) =
\sum_{k=1}^\infty f_k(x) .
\end{equation*}

\medskip

Pertanyaannya:
Jika semua $f_n$ kontinu, apakah $f$ kontinu?  Diferensiabel?
Dapat diintegralkan?  Apa turunan atau integral dari $f$?
Sebagai contoh, untuk kekontinuan limit titik demi titik dari suatu barisan
fungsi
$\{ f_n \}_{n=1}^\infty$, kita menanyakan apakah
\begin{equation*}
\lim_{x\to x_0} \lim_{n\to\infty} f_n(x)
\overset{?}{=}
\lim_{n\to\infty} \lim_{x\to x_0} f_n(x) .
\end{equation*}
Secara apriori, kita bahkan tidak tahu apakah kedua ruas tersebut ada,
apalagi apakah keduanya sama.

\begin{example}
Fungsi $f_n \colon \R \to \R$,
\begin{equation*}
f_n(x) \coloneqq \frac{1}{1+nx^2},
\end{equation*}
kontinu dan konvergen titik demi titik ke fungsi yang tidak kontinu
\begin{equation*}
f(x) \coloneqq 
\begin{cases}
1 & \text{jika } x=0, \\
0 & \text{selainnya.}
\end{cases}
\end{equation*}
\end{example}

Jadi, konvergensi titik demi titik tidak cukup untuk mempertahankan kekontinuan
(bahkan keterbatasan). Untuk itu, kita memerlukan konvergensi seragam.
Misalkan $f_n \colon X \to Y$ fungsi-fungsi. Maka
$\{f_n\}_{n=1}^\infty$ \emph{\myindex{konvergen seragam}}\index{konvergensi seragam}
ke $f$ jika
untuk setiap $\epsilon > 0$, terdapat suatu $M$ sedemikian sehingga
untuk semua $n \geq M$ dan semua $x \in X$, berlaku
\begin{equation*}
d\bigl(f_n(x),f(x)\bigr) < \epsilon .
\end{equation*}

Deret $\sum_{n=1}^\infty f_n$ dari fungsi-fungsi bernilai kompleks konvergen seragam
ke $f$ jika barisan
jumlah parsialnya konvergen seragam, yaitu, jika untuk setiap $\epsilon > 0$,
terdapat suatu $M$ sedemikian sehingga
untuk semua $n \geq M$ dan semua $x \in X$,
\begin{equation*}
\abs{\left(\sum_{k=1}^n f_k(x)\right)-f(x)} < \epsilon .
\end{equation*}

Sifat paling sederhana yang dipertahankan oleh konvergensi seragam adalah
keterbatasan. Pembuktian proposisi berikut kami serahkan sebagai
latihan. Pembuktiannya hampir identik dengan pembuktian untuk fungsi bernilai real.

\begin{prop} \label{prop:uniformconvbounded}
Misalkan $X$ suatu himpunan dan $(Y,d)$ suatu ruang metrik.
Jika $f_n \colon X \to Y$ adalah fungsi-fungsi terbatas dan konvergen seragam ke $f
\colon X \to Y$, maka $f$ terbatas.
\end{prop}

Jika $X$ suatu himpunan dan $(Y,d)$ suatu ruang metrik, maka suatu barisan $f_n \colon X
\to Y$ disebut \emph{\myindex{Cauchy seragam}} jika untuk setiap $\epsilon > 0$, terdapat suatu $M$ sedemikian sehingga
untuk semua $n, m \geq M$ dan semua $x \in X$, berlaku
\begin{equation*}
d\bigl(f_n(x),f_m(x)\bigr) < \epsilon .
\end{equation*}
Gagasan ini sama seperti untuk fungsi bernilai real.
Pembuktian proposisi berikut
pada dasarnya juga sama seperti pada kasus tersebut dan
diserahkan sebagai latihan.

\begin{prop} \label{prop:unifcauchymetric}
Misalkan $X$ suatu himpunan, $(Y,d)$ suatu ruang metrik, dan
$f_n \colon X \to Y$ fungsi-fungsi.
Jika $\{ f_n \}_{n=1}^\infty$ konvergen seragam,
maka $\{f_n\}_{n=1}^\infty$ bersifat Cauchy seragam. Sebaliknya, jika
$\{f_n\}_{n=1}^\infty$ bersifat Cauchy seragam dan $(Y,d)$ lengkap secara Cauchy,
maka $\{f_n\}_{n=1}^\infty$ konvergen seragam.
\end{prop}

Untuk $f \colon X \to \C$, kita tuliskan\glsadd{not:uniformnorm}
\begin{equation*}
\snorm{f}_X \coloneqq \sup_{x \in X} \babs{f(x)} .
\end{equation*}
Kita menyebut $\snorm{\cdot}_X$
sebagai \emph{\myindex{norma supremum}} atau \emph{\myindex{norma seragam}},
dan subskripnya menunjukkan himpunan tempat supremum diambil.
Maka suatu barisan fungsi
$f_n \colon X \to \C$ konvergen seragam ke $f \colon X \to \C$ jika dan hanya jika
\begin{equation*}
\lim_{n\to \infty} \snorm{f_n-f}_X = 0 .
\end{equation*}

Norma supremum memenuhi ketaksamaan segitiga: Untuk setiap $x \in X$,
\begin{equation*}
\babs{f(x)+g(x)} \leq
\babs{f(x)}+\babs{g(x)} \leq
\snorm{f}_X+\snorm{g}_X .
\end{equation*}
Ambil supremum pada ruas kiri untuk memperoleh
\begin{equation*}
\snorm{f+g}_X \leq
\snorm{f}_X+\snorm{g}_X .
\end{equation*}

Untuk ruang metrik kompak $X$,
norma seragam merupakan norma pada ruang vektor $C(X,\C)$.
Kami menyerahkan pembuktiannya sebagai latihan.
Meskipun kita tidak akan memerlukannya, $C(X,\C)$ sebenarnya merupakan ruang
vektor kompleks, yaitu, dalam definisi ruang vektor kita dapat mengganti
$\R$ dengan $\C$.
Konvergensi dalam ruang metrik $C(X,\C)$ adalah
konvergensi seragam.

Kita akan mempelajari beberapa jenis deret fungsi, dan
suatu uji yang berguna untuk konvergensi seragam suatu deret adalah
\emph{\myindex{uji-$M$ Weierstrass}}.

\begin{thm}[\myindex{uji-$M$ Weierstrass}] \label{thm:weiermtest}
Misalkan $X$ suatu himpunan.
Andaikan $f_n \colon X \to \C$ adalah fungsi-fungsi dan $M_n > 0$ adalah bilangan-bilangan sedemikian
sehingga
\begin{equation*}
\babs{f_n(x)}\leq M_n \quad \text{untuk semua } x \in X,
\qquad \text{dan} \qquad
\sum_{n=1}^\infty M_n
\quad \text{konvergen}.
\end{equation*}
Maka
\begin{equation*}
\sum_{n=1}^\infty f_n(x)
\quad \text{konvergen seragam}.
\end{equation*}
\end{thm}

Cara lain untuk menyatakan teorema tersebut adalah dengan mengatakan bahwa jika
$\sum_{n=1}^\infty \snorm{f_n}_X$ konvergen, maka $\sum_{n=1}^\infty f_n$ konvergen seragam.
Perhatikan bahwa kebalikan teorema ini tidak benar.
Dengan menerapkan teorema tersebut pada $\sum_{n=1}^\infty \babs{f_n(x)}$, kita
melihat bahwa deret ini juga konvergen seragam. Jadi, deret tersebut konvergen absolut
sekaligus seragam.

\begin{proof}
Andaikan $\sum_{n=1}^\infty M_n$ konvergen. Ambil $\epsilon > 0$.
Jumlah-jumlah parsial deret $\sum_{n=1}^\infty M_n$ bersifat Cauchy, sehingga
terdapat suatu $N$ sedemikian sehingga untuk semua $m, n \geq N$ dengan $m > n$, berlaku
\begin{equation*}
\sum_{k=n+1}^m M_k < \epsilon .
\end{equation*}
Kita menaksir selisih Cauchy antara jumlah-jumlah parsial
fungsi-fungsi tersebut:
\begin{equation*}
\abs{\sum_{k=n+1}^m f_k(x)} \leq
\sum_{k=n+1}^m \babs{f_k(x)} \leq
\sum_{k=n+1}^m M_k < \epsilon .
\end{equation*}
Deret tersebut konvergen berdasarkan \propref{prop:cachysercomplex}.
Konvergensinya seragam karena $N$ tidak bergantung pada $x$.
Sesungguhnya, untuk semua $n \geq N$,
\begin{equation*}
\abs{\sum_{k=1}^\infty f_k(x) - \sum_{k=1}^n f_k(x)} \leq
\abs{\sum_{k=n+1}^\infty f_k(x)} \leq \epsilon .
\qedhere
\end{equation*}
\end{proof}

\begin{example} \label{example:sinnsqfourier}
Deret
\begin{equation*}
\sum_{n=1}^\infty \frac{\sin(nx)}{n^2}
\end{equation*}
konvergen seragam pada $\R$. Lihat \figureref{fig:fouriersern2}.
Deret ini merupakan deret Fourier, dan
kita akan melihat lebih banyak deret semacam ini pada bagian selanjutnya. Pembuktian:
Deret tersebut konvergen seragam karena
$\sum_{n=1}^\infty \frac{1}{n^2}$
konvergen dan
\begin{equation*}
\abs{\frac{\sin(nx)}{n^2}} \leq 
\frac{1}{n^2} .
\end{equation*}
\end{example}

\begin{myfigureht}
\myincludegraphics{fouriersern2}{%
Grafik suatu fungsi yang digambarkan dengan garis tebal, berbentuk
gelombang dengan lereng turun yang landai dan lereng naik yang curam.
Beberapa aproksimasi diperlihatkan dalam berbagai nuansa abu-abu;
aproksimasi yang lebih gelap (lebih awal) lebih menyerupai gelombang sinus biasa,
sedangkan yang lebih terang lebih menyerupai garis tebal gelap.}
\caption{Grafik
$\sum_{n=1}^\infty \frac{\sin(n x)}{n^2}$ yang mencakup
8 jumlah parsial pertamanya dalam berbagai nuansa abu-abu.\label{fig:fouriersern2}}
\end{myfigureht}

\begin{example}
Deret
\begin{equation*}
\sum_{n=0}^\infty \frac{x^n}{n!} 
\end{equation*}
konvergen seragam pada setiap interval terbatas.
Deret ini merupakan deret pangkat yang akan segera kita pelajari.
Pembuktian: Ambil interval $[-r,r] \subset \R$ (setiap interval terbatas
termuat dalam suatu $[-r,r]$).
Deret $\sum_{n=0}^\infty \frac{r^n}{n!}$ konvergen berdasarkan uji rasio,
sehingga $\sum_{n=0}^\infty \frac{x^n}{n!}$ konvergen seragam pada $[-r,r]$ karena
\begin{equation*}
\abs{\frac{x^n}{n!} } \leq 
\frac{r^n}{n!} .
\end{equation*}
\end{example}

Sekarang kita ingin mengetahui sesuatu tentang limit tersebut. Sebagai contoh,
apakah limit itu kontinu?


\begin{prop} \label{prop:uniformswitch}
Misalkan $(X,d_X)$ dan $(Y,d_Y)$ ruang-ruang metrik, dan andaikan $(Y,d_Y)$
lengkap secara Cauchy.
Andaikan $f_n \colon X \to Y$ konvergen seragam ke
 $f \colon X \to Y$.
Misalkan $\{ x_k \}_{k=1}^\infty$ suatu barisan dalam $X$ dan $x \coloneqq \lim_{k \to \infty} x_k$. Andaikan
\begin{equation*}
a_n \coloneqq \lim_{k \to \infty} f_n(x_k)
\end{equation*}
ada untuk semua $n$. Maka
$\{a_n\}_{n=1}^\infty$ konvergen dan
\begin{equation*}
\lim_{k \to \infty} f(x_k) = \lim_{n\to\infty} a_n .
\end{equation*}
\end{prop}

Dengan kata lain,
\begin{equation*}
\lim_{k \to \infty} \lim_{n\to\infty} f_n(x_k) =
\lim_{n \to \infty} \lim_{k\to\infty} f_n(x_k) .
\end{equation*}

\begin{proof}
Pertama, kita tunjukkan bahwa $\{ a_n \}_{n=1}^\infty$ konvergen. Karena
$\{ f_n \}_{n=1}^\infty$ konvergen seragam, barisan ini bersifat Cauchy seragam.
Ambil $\epsilon > 0$. Terdapat
suatu $M$ sedemikian sehingga untuk semua $m,n \geq M$, berlaku
\begin{equation*}
d_Y\bigl(f_n(x_k),f_m(x_k)\bigr) < \epsilon \qquad \text{untuk semua } k .
\end{equation*}
Perhatikan bahwa
$d_Y(a_n,a_m) \leq
d_Y\bigl(a_n,f_n(x_k)\bigr) +
d_Y\bigl(f_n(x_k),f_m(x_k)\bigr) +
d_Y\bigl(f_m(x_k),a_m\bigr)$. Dengan mengambil limit ketika $k \to \infty$, kita memperoleh
\begin{equation*}
d_Y(a_n,a_m) \leq \epsilon .
\end{equation*}
Karena itu, $\{a_n\}_{n=1}^\infty$ bersifat Cauchy dan konvergen karena $Y$ lengkap. Tuliskan
$a \coloneqq \lim_{n \to \infty} a_n$.

Pilih suatu $\ell \in \N$ sedemikian sehingga
\begin{equation*}
d_Y\bigl(f_\ell(p),f(p)\bigr) < \nicefrac{\epsilon}{3}
\end{equation*}
untuk semua $p \in X$. Pastikan pula $\ell$ cukup besar
sehingga
\begin{equation*}
d_Y(a_\ell,a) < \nicefrac{\epsilon}{3}  .
\end{equation*}
Pilih suatu $N \in \N$ sedemikian sehingga untuk $m \geq N$,
\begin{equation*}
d_Y\bigl(f_\ell(x_m),a_\ell\bigr) < \nicefrac{\epsilon}{3}  .
\end{equation*}
Maka untuk
$m \geq N$,
\begin{equation*}
d_Y\bigl(f(x_m),a\bigr)
\leq
d_Y\bigl(f(x_m),f_\ell(x_m)\bigr)
+
d_Y\bigl(f_\ell(x_m),a_\ell\bigr)
+
d_Y\bigl(a_\ell,a\bigr)
<
\nicefrac{\epsilon}{3} +
\nicefrac{\epsilon}{3} +
\nicefrac{\epsilon}{3} = \epsilon . \qedhere
\end{equation*}
\end{proof}

Kita memperoleh korolari langsung mengenai kekontinuan.
Jika semua $f_n$ kontinu, maka $a_n = f_n(x)$.
Karena itu,
$\{ a_n \}_{n=1}^\infty$ secara otomatis konvergen ke
$f(x)$ berdasarkan hipotesis, sehingga kita tidak memerlukan kelengkapan $Y$.

\begin{cor} \label{cor:metricuniformcontinuous}
Misalkan $X$ dan $Y$ ruang-ruang metrik.
Jika $f_n \colon X \to Y$ fungsi-fungsi kontinu
sedemikian sehingga
$\{ f_n \}_{n=1}^\infty$ konvergen seragam ke $f \colon X \to Y$,
maka $f$ kontinu.
\end{cor}

Kebalikannya tidak benar. Fakta bahwa limitnya kontinu tidak berarti
bahwa konvergensinya seragam. Sebagai contoh:
$f_n \colon (0,1) \to \R$ yang didefinisikan oleh $f_n(x) \coloneqq x^n$ konvergen ke
fungsi nol, tetapi tidak secara seragam. Namun, jika kita menambahkan syarat-syarat tambahan
pada barisan tersebut, kita dapat memperoleh kebalikan parsial seperti teorema Dini,
\volIref{lihat \exerciseref*{vI-exercise:dinisthm} dari jilid I}{lihat \exerciseref{exercise:dinisthm}}.

Dalam \exerciseref{exercise:CXCnormedspace}, pembaca diminta untuk membuktikan
bahwa untuk $X$ kompak, $C(X,\C)$ merupakan ruang vektor bernorma
dengan norma seragam dan dengan demikian merupakan ruang metrik. Kita baru saja menunjukkan bahwa
$C(X,\C)$ lengkap secara Cauchy: \propref{prop:unifcauchymetric} mengatakan bahwa suatu barisan Cauchy
dalam $C(X,\C)$ konvergen seragam ke suatu fungsi,
dan \corref{cor:metricuniformcontinuous} menunjukkan bahwa limitnya
kontinu dan dengan demikian merupakan anggota $C(X,\C)$.

\begin{cor}
Misalkan $(X,d)$ suatu ruang metrik kompak.
Maka $C(X,\C)$ merupakan ruang metrik yang lengkap secara Cauchy.
\end{cor}

\begin{example}
Berdasarkan \exampleref{example:sinnsqfourier},
deret Fourier
\begin{equation*}
\sum_{n=1}^\infty \frac{\sin(nx)}{n^2}
\end{equation*}
konvergen seragam dan karenanya kontinu berdasarkan \corref{cor:metricuniformcontinuous} (seperti yang tampak
pada \figureref{fig:fouriersern2}).
\end{example}

\subsection{Integrasi}

\begin{prop} \label{prop:complexlimitswapintegral}
Misalkan $f_n \colon [a,b] \to \C$
terintegralkan secara Riemann dan andaikan $\{ f_n \}_{n=1}^\infty$ konvergen
seragam ke $f \colon [a,b] \to \C$. Maka $f$ terintegralkan secara Riemann
dan
\begin{equation*}
\int_a^b f = \lim_{n\to \infty} \int_a^b f_n .
\end{equation*}
\end{prop}

Karena integral suatu fungsi bernilai kompleks hanyalah integral
bagian real dan imajinernya secara terpisah,
pembuktiannya langsung mengikuti hasil-hasil dalam
\volIref{\chapterref*{vI-fs:chapter} dari jilid~I}{\chapterref{fs:chapter}}.
Rinciannya kita serahkan sebagai latihan.

\begin{cor}
\pagebreak[2]
Misalkan $f_n \colon [a,b] \to \C$
terintegralkan secara Riemann dan andaikan bahwa
\begin{equation*}
\sum_{n=1}^\infty f_n(x)
\end{equation*}
konvergen seragam. Maka deret tersebut terintegralkan secara Riemann pada $[a,b]$
dan
\begin{equation*}
\int_a^b \sum_{n=1}^\infty f_n(x) \,dx
=
\sum_{n=1}^\infty \int_a^b f_n(x) \,dx
\end{equation*}
\end{cor}

\begin{example}
Mari kita tunjukkan cara mengintegralkan suatu deret Fourier.
\begin{equation*}
\int_{0}^x \sum_{n=1}^\infty \frac{\cos(nt)}{n^2} \,dt
=
\sum_{n=1}^\infty \int_{0}^x \frac{\cos(nt)}{n^2}\,dt
=
\sum_{n=1}^\infty \frac{\sin(nx)}{n^3}
\end{equation*}
Pertukaran integral dan jumlah dimungkinkan oleh konvergensi seragam,
yang sebelumnya telah kita buktikan menggunakan uji-$M$ Weierstrass
(\thmref{thm:weiermtest}).
\end{example}

Perhatikan bahwa kita dapat mempertukarkan integral dan limit di bawah hipotesis yang jauh lebih lemah,
tetapi untuk itu kita memerlukan integral yang lebih kuat daripada integral Riemann,
seperti integral Lebesgue.

\subsection{Diferensiasi}

Ingat bahwa suatu fungsi bernilai kompleks
$f \colon [a,b] \to \C$, dengan $f(x) = u(x)+i\,v(x)$,
diferensiabel jika $u$ dan $v$ diferensiabel
dan turunannya adalah
\begin{equation*}
f'(x) = u'(x)+i\,v'(x) .
\end{equation*}

Pembuktian teorema berikut dilakukan dengan menerapkan teorema yang bersesuaian untuk
fungsi real pada $u$ dan $v$, dan diserahkan sebagai latihan.

\begin{thm} \label{thm:dersconvergecomplex}
Misalkan $I \subset \R$ suatu interval terbatas dan misalkan
$f_n \colon I \to \C$ fungsi-fungsi yang diferensiabel secara kontinu.
Andaikan $\{ f_n' \}_{n=1}^\infty$ konvergen seragam ke $g \colon I \to \C$,
dan andaikan $\{ f_n(c) \}_{n=1}^\infty$ merupakan suatu
barisan konvergen untuk suatu $c \in I$.  Maka $\{ f_n \}_{n=1}^\infty$ konvergen seragam ke 
suatu fungsi $f \colon I \to \C$ yang diferensiabel secara kontinu, dan $f' = g$.
\end{thm}

Konvergensi seragam dari fungsi-fungsi itu sendiri tidaklah cukup.
Turunannya mungkin tidak konvergen dan fungsi limitnya bahkan tidak harus
diferensiabel.
Dalam \sectionref{sec:stoneweier}, kita akan membuktikan bahwa
fungsi-fungsi kontinu merupakan limit seragam dari polinom, tetapi seperti yang ditunjukkan oleh
contoh berikut, suatu fungsi kontinu tidak harus diferensiabel
di titik mana pun.

\begin{example}
Terdapat fungsi kontinu yang tidak diferensiabel di mana pun.
Fungsi semacam itu sering disebut
\emph{fungsi Weierstrass}\index{fungsi Weierstrass},
meskipun fungsi
khusus ini, yang pada hakikatnya berasal dari
Takagi\footnote{\href{https://en.wikipedia.org/wiki/Teiji_Takagi}{Teiji
Takagi} (1875--1960) adalah seorang matematikawan Jepang.}, merupakan contoh yang berbeda dari yang diberikan Weierstrass.
Definisikan
\begin{equation*}
\varphi(x) \coloneqq \sabs{x} \qquad \text{untuk } x \in [-1,1] .
\end{equation*}
Perluas $\varphi$ ke seluruh $\R$ dengan menjadikannya
2-periodik:
Tetapkan bahwa
$\varphi(x) = \varphi(x+2)$.  Fungsi $\varphi \colon \R \to \R$
kontinu; bahkan, $\babs{\varphi(x)-\varphi(y)} \leq \sabs{x-y}$ (mengapa?).
Lihat \figureref{fig:triangwave}.
\begin{myfigureht}
\myincludegraphics{triangwave}{%
Suatu fungsi yang grafiknya berupa deretan segitiga; yakni,
fungsi tersebut bernilai 0 pada semua bilangan bulat genap, dan bernilai 1 pada semua bilangan bulat
ganjil.  Di antaranya, grafik tersebut berupa garis lurus.}
\caption{Fungsi 2-periodik $\varphi$.\label{fig:triangwave}}
\end{myfigureht}

Karena $\sum_{n=0}^\infty {\left(\frac{3}{4}\right)}^n$ konvergen dan $\babs{\varphi(x)} \leq
1$ untuk semua $x$, menurut uji-$M$
(\thmref{thm:weiermtest}),
\begin{equation*}
f(x) \coloneqq \sum_{n=0}^\infty 
{\left(\frac{3}{4}\right)}^n \varphi(4^n x)
\end{equation*}
konvergen seragam sehingga kontinu.
Lihat \figureref{fig:nowherediff}.

\begin{myfigureht}
\myincludegraphics{nowherediff}{%
Grafik suatu fungsi yang sangat liar dengan banyak 
puncak tinggi yang runcing dan titik minimum yang rendah, seakan-akan digambar dengan menggerakkan pena secara sangat
cepat naik turun.}
\caption{Grafik fungsi $f$ yang tidak diferensiabel di mana pun.\label{fig:nowherediff}}
\end{myfigureht}

Kita mengklaim bahwa $f \colon
\R \to \R$ tidak diferensiabel di mana pun.
Ambil sembarang $x$, dan akan kita tunjukkan bahwa $f$ tidak diferensiabel di $x$.
Definisikan
\begin{equation*}
\delta_m \coloneqq \pm \frac{1}{2} 4^{-m} ,
\avoidbreak
\end{equation*}
dengan tanda dipilih sedemikian sehingga tidak ada bilangan bulat
di antara $4^m x$ dan $4^m(x+\delta_m) = 4^m x \pm \frac{1}{2}$.

Kita akan meninjau hasil bagi selisih
\begin{equation*}
\frac{f(x+\delta_m)-f(x)}{\delta_m}
=
\sum_{n=0}^\infty 
{\left(\frac{3}{4}\right)}^n
\frac{\varphi\bigl(4^n(x+\delta_m)\bigr)-\varphi(4^nx)}{\delta_m} .
\end{equation*}
Untuk sementara, tetapkan $m$.  Tinjau ekspresi di dalam deret:
\begin{equation*}
\gamma_{n} \coloneqq
\frac{\varphi\bigl(4^n(x+\delta_m)\bigr)-\varphi(4^nx)}{\delta_m} .
\end{equation*}
Jika $n > m$, maka $4^n\delta_m$ merupakan bilangan bulat genap.  Karena $\varphi$
2-periodik, diperoleh $\gamma_n = 0$.

Karena tidak ada bilangan bulat di antara 
$4^m(x+\delta_m) = 4^m x\pm\nicefrac{1}{2}$ dan $4^m x$, maka pada interval ini
$\varphi(t) = \pm t + \ell$ untuk suatu bilangan bulat $\ell$.
Secara khusus,
$\abs{\varphi\bigl(4^m(x+ \delta_m)\bigr)-\varphi(4^mx)} =
\abs{4^mx\pm\nicefrac{1}{2}-4^mx} = \nicefrac{1}{2}$.  Oleh karena itu,
\begin{equation*}
\sabs{\gamma_m} =
\abs{
\frac{\varphi\bigl(4^m(x+\delta_m)\bigr)-\varphi(4^mx)}{\pm (\nicefrac{1}{2}) 4^{-m}}
}
= 4^m .
\end{equation*}
Demikian pula, andaikan $n < m$.  Karena $\babs{\varphi(s) -\varphi(t)} \leq
\sabs{s-t}$,
\begin{equation*}
\sabs{\gamma_n} =
\abs{\frac{\varphi\bigl(4^nx\pm(\nicefrac{1}{2})4^{n-m}\bigr)-\varphi(4^nx)}{\pm
(\nicefrac{1}{2}) 4^{-m}}}
\leq
\abs{\frac{\pm(\nicefrac{1}{2})4^{n-m}}{\pm (\nicefrac{1}{2}) 4^{-m}}} = 4^n
.
\end{equation*}

Dengan demikian,
\begin{equation*}
\begin{split}
\abs{
\frac{f(x+\delta_m)-f(x)}{\delta_m}
}
%& =
%\abs{
%\sum_{n=0}^\infty 
%{\left(\frac{3}{4}\right)}^n
%\frac{\varphi\bigl(4^n(x+\delta_m)\bigr)-\varphi(4^nx)}{\delta_m}
%}
=
\abs{
\sum_{n=0}^\infty 
{\left(\frac{3}{4}\right)}^n
\gamma_n
}
& =
\abs{
\sum_{n=0}^m 
{\left(\frac{3}{4}\right)}^n
\gamma_n
}
\\
& \geq
\abs{
{\left(\frac{3}{4}\right)}^m
\gamma_m}
-
\abs{
\sum_{n=0}^{m-1} 
{\left(\frac{3}{4}\right)}^n
\gamma_n
}
\\
& \geq
3^m
-
\sum_{n=0}^{m-1} 
3^n
=
3^m
-
\frac{3^{m}-1}{3-1}
=
\frac{3^m +1}{2} .
\end{split}
\end{equation*}
Ketika $m \to \infty$, kita memperoleh
$\delta_m \to 0$, tetapi $\frac{3^m+1}{2}$
menuju tak hingga.  Jadi $f$ tidak mungkin diferensiabel di~$x$.
\end{example}

\subsection{Latihan}

\begin{exercise}
Buktikan \propref{prop:uniformconvbounded}.
\end{exercise}

\begin{exercise}
Buktikan \propref{prop:unifcauchymetric}.
\end{exercise}

\begin{exercise} \label{exercise:CXCnormedspace}
Misalkan $(X,d)$ adalah ruang metrik kompak.
Buktikan bahwa norma seragam $\snorm{\cdot}_X$ merupakan norma pada ruang vektor
fungsi-fungsi kontinu bernilai kompleks $C(X,\C)$.
\end{exercise}

\begin{exercise}
\pagebreak[2]
\leavevmode
\begin{enumerate}[a)]
\item
Buktikan bahwa
$f_n(x) \coloneqq 2^{-n} \sin(2^n x)$
konvergen seragam ke nol, tetapi terdapat himpunan padat $D \subset \R$
sedemikian sehingga $\lim_{n\to\infty} f_n'(x) = 1$ untuk setiap $x \in D$.
\item
Buktikan bahwa
$\sum_{n=1}^\infty 2^{-n} \sin(2^n x)$
konvergen seragam ke suatu fungsi kontinu,
dan terdapat himpunan padat $D \subset \R$
sehingga turunan jumlah-jumlah parsial tersebut tidak konvergen di setiap titik himpunan itu.
\end{enumerate}
\end{exercise}

\begin{exercise}
Buktikan bahwa $\snorm{f}_{C^1} \coloneqq \snorm{f}_{[a,b]}+\snorm{f'}_{[a,b]}$
merupakan norma pada ruang vektor fungsi-fungsi bernilai kompleks yang
diferensiabel secara kontinu $C^1\bigl([a,b],\C\bigr)$.
\end{exercise}

\begin{exercise}
Buktikan \thmref{thm:dersconvergecomplex}.
\end{exercise}

\begin{exercise}
Buktikan \propref{prop:complexlimitswapintegral} dengan mereduksinya ke
hasil untuk kasus real.
\end{exercise}

\begin{exercise}
Kerjakan rincian contoh penyangkal berikut untuk kebalikan
uji-$M$ Weierstrass (\thmref{thm:weiermtest}).  Definisikan
$f_n \colon [0,1] \to \R$ melalui
\begin{equation*}
f_n(x) \coloneqq
\begin{cases}
\frac{1}{n} & \text{jika } \frac{1}{n+1} < x < \frac{1}{n},\\
0	    & \text{selainnya.}
\end{cases}
\end{equation*}
Buktikan bahwa $\sum_{n=1}^\infty f_n$ konvergen seragam, tetapi
$\sum_{n=1}^\infty \snorm{f_n}_{[0,1]}$
tidak konvergen.
\end{exercise}

\begin{exercise}
Misalkan $f_n \colon [0,1] \to \R$ adalah fungsi-fungsi yang monoton naik dan
misalkan $\sum_{n=1}^\infty f_n$ konvergen titik demi titik.  Buktikan bahwa
$\sum_{n=1}^\infty f_n$
konvergen seragam.
\end{exercise}

\begin{exercise}
Buktikan bahwa
\begin{equation*}
\sum_{n=1}^\infty e^{-nx}
\end{equation*}
konvergen untuk setiap $x > 0$ ke suatu fungsi yang diferensiabel.
\end{exercise}

%%%%%%%%%%%%%%%%%%%%%%%%%%%%%%%%%%%%%%%%%%%%%%%%%%%%%%%%%%%%%%%%%%%%%%%%%%%%%%

\sectionnewpage
\section{Deret pangkat dan fungsi analitik}
\label{sec:analfuncs}

%mbxINTROSUBSECTION

\sectionnotes{2--3 kuliah}

\subsection{Fungsi analitik}

Suatu deret pangkat (kompleks) adalah deret berbentuk
\begin{equation*}
\sum_{n=0}^\infty c_n {(z-a)}^n
\end{equation*}
untuk $c_n, z, a \in \C$.  Kita mengatakan deret tersebut
\emph{konvergen}\index{konvergen!deret pangkat} jika deret itu konvergen untuk
suatu $z \neq a$.

Misalkan $U \subset \C$ merupakan himpunan terbuka dan
$f \colon U \to \C$ suatu fungsi.
Andaikan untuk setiap $a \in U$ terdapat $\rho > 0$ dan suatu deret
pangkat yang konvergen ke fungsi tersebut,
\begin{equation*}
f(z) = \sum_{n=0}^\infty c_n {(z-a)}^n
\end{equation*}
untuk semua $z \in B(a,\rho)$.
Maka kita mengatakan bahwa $f$ adalah fungsi \emph{\myindex{analitik}}.
Serupa dengan itu, untuk interval $(a,b) \subset \R$, kita mengatakan bahwa $f \colon (a,b) \to \C$
bersifat analitik atau dapat pula disebut \emph{\myindex{analitik real}}
jika untuk setiap titik $c \in
(a,b)$ terdapat deret pangkat di sekitar $c$ yang konvergen pada suatu
$(c-\rho,c+\rho)$
untuk suatu $\rho > 0$.
Karena kita kadang-kadang akan membahas deret pangkat real dan kadang-kadang
deret pangkat kompleks, kita akan memakai $z$ untuk menyatakan bilangan kompleks dan $x$ bilangan real.
Kita akan selalu menyebutkan kasus mana yang sedang dibahas.

Suatu fungsi analitik memiliki pengembangan yang berbeda di sekitar titik-titik yang berbeda.
Selain itu, konvergensi tidak otomatis terjadi pada seluruh
domain fungsi.  Sebagai contoh, jika $\sabs{z} < 1$, maka
\begin{equation*}
\frac{1}{1-z} = \sum_{n=0}^\infty z^n .
\end{equation*}
Walaupun ruas kiri terdefinisi untuk setiap $z \neq 1$, ruas kanan
ternyata hanya konvergen jika $\sabs{z} < 1$.  Lihat grafik
sebagian kecil dari $\frac{1}{1-z}$ pada \figureref{fig:1over1mz}.
Kita tidak dapat menggambar grafik
fungsi itu sendiri; kita hanya dapat menggambar bagian real atau bagian imajinernya karena
dimensi alam semesta kita tidak mencukupi.

\begin{myfigureht}
\myincludepdft{real_imag_1over1mz}{%
Dua grafik permukaan 3D di atas bidang xy.  Pada kedua grafik, suatu garis vertikal
ditandai di atas titik (1,0) pada bidang xy.  Grafik-grafik tersebut ditampilkan pada
petak kecil di sekitar titik asal dan berakhir sebelum mencapai titik (1,0).
Grafik kiri lebih datar jauh dari garis vertikal dan menanjak
curam ketika mendekati garis vertikal.  Grafik kanan serupa,
tetapi menanjak curam ketika mendekati garis vertikal dari
daerah tempat y bernilai negatif dan menurun curam ketika mendekati
garis vertikal dari daerah tempat y bernilai positif.}
\caption{Grafik bagian real dan bagian imajiner dari $z=x+iy \mapsto \frac{1}{1-z}$
pada persegi $[-0.8,0.8]^2$.  Singularitas di $z=1$ ditandai dengan
garis vertikal putus-putus.\label{fig:1over1mz}}
\end{myfigureht}

\subsection{Konvergensi deret pangkat}

Kita telah membuktikan beberapa hasil untuk deret pangkat dengan peubah real dalam
\volIref{\sectionref*{vI-sec:moreonseries} pada Jilid I}{\sectionref{sec:moreonseries}}.   Sebagian besar sifat
konvergensi deret pangkat berkaitan dengan deret
$\sum_{k=0}^\infty \sabs{c_k} \, \sabs{z-a}^k$,
sehingga kita sebenarnya telah membuktikan banyak hasil mengenai deret pangkat
kompleks.  Khususnya, kita telah menghitung apa yang disebut
jari-jari konvergensi suatu deret pangkat.

\begin{prop}
\pagebreak[2]
Misalkan $\sum_{n=0}^\infty c_n {(z-a)}^n$ suatu deret pangkat.
Terdapat $\rho \in [0,\infty]$ sedemikian sehingga
\begin{enumerate}[(i)]
\item Jika $\rho = 0$, maka deret tersebut divergen untuk setiap $z \neq a$.
\item Jika $\rho = \infty$, maka deret tersebut konvergen mutlak untuk semua $z \in \C$.
\item Jika $0 < \rho < \infty$, maka deret tersebut
konvergen mutlak pada $B(a,\rho)$,
dan divergen ketika $\sabs{z-a} > \rho$.
\end{enumerate}
Lebih lanjut, jika $0 < r < \rho$, maka
deret tersebut konvergen secara seragam pada bola tertutup $C(a,r)$.
\end{prop}

Bilangan $\rho$ disebut \emph{\myindex{jari-jari konvergensi}}.
Lihat \figureref{fig:radiusconvcomplex}.
Jari-jari konvergensi menentukan sebuah cakram berpusat di $a$ tempat deret tersebut konvergen.  Suatu deret pangkat
dikatakan konvergen jika $\rho > 0$.
\begin{myfigureht}
\myincludepdft{radiusconvcomplex}{%
Sebuah lingkaran bertitik-titik berpusat di a dengan jari-jari rho.
Bagian dalam lingkaran diarsir dan diberi label "deret konvergen".
Bagian luar lingkaran diberi label "deret tidak konvergen".}
\caption{Jari-jari konvergensi.\label{fig:radiusconvcomplex}}
\end{myfigureht}

\begin{proof}
Kita menggunakan versi real dari proposisi ini,
\volIref{\propref*{vI-prop:powerserrealradius} pada Jilid I}{\propref{prop:powerserrealradius}}.
Misalkan
\begin{equation*}
R \coloneqq \limsup_{n\to\infty} \sqrt[n]{\sabs{c_n}} .
\end{equation*}
Jika $R = 0$, maka
$\sum_{n=0}^\infty \sabs{c_n} \, \sabs{z-a}^n$ konvergen untuk semua $z$.
Jika $R = \infty$, maka
$\sum_{n=0}^\infty \sabs{c_n} \, \sabs{z-a}^n$ hanya konvergen di $z=a$.
Jika tidak, misalkan $\rho \coloneqq \nicefrac{1}{R}$.
Maka
$\sum_{n=0}^\infty \sabs{c_n} \, \sabs{z-a}^n$ konvergen ketika
$\sabs{z-a} < \rho$, dan divergen (bahkan suku-suku deretnya
tidak menuju nol) ketika $\sabs{z-a} > \rho$.

Untuk membuktikan pernyataan \myquote{Lebih lanjut,} andaikan
$0 < r < \rho$ dan $z \in C(a,r)$.  Perhatikan
jumlah parsial
\begin{equation*}
\abs{\sum_{n=0}^k c_n {(z-a)}^n}
\leq
\sum_{n=0}^k \sabs{c_n} \sabs{z-a}^n
\leq
\sum_{n=0}^k \sabs{c_n} r^n . \qedhere
\end{equation*}
\end{proof}

Jika $\sum_{n=0}^\infty c_n {(z-a)}^n$ konvergen untuk suatu $z$, maka
\begin{equation*}
\sum_{n=0}^\infty c_n {(w-a)}^n
\end{equation*}
konvergen mutlak setiap kali $\sabs{w-a} < \sabs{z-a}$.
Sebaliknya, jika deret tersebut divergen di $z$, maka deret itu juga harus divergen di $w$
setiap kali $\sabs{w-a} > \sabs{z-a}$.
Jadi, untuk menunjukkan
bahwa jari-jari konvergensinya sekurang-kurangnya sebesar suatu bilangan tertentu, kita hanya perlu
menunjukkan konvergensi di suatu titik dengan metode apa pun yang kita ketahui.

\begin{example}
Kita daftarkan beberapa deret yang telah kita kenal:
\begin{align*}
& &
& \sum_{n=0}^\infty z^n
& & \text{memiliki jari-jari konvergensi } 1.
& &
\\
& &
& \sum_{n=0}^\infty \frac{1}{n!} z^n
& & \text{memiliki jari-jari konvergensi } \infty.
& &
\\
& &
& \sum_{n=0}^\infty n^n z^n
& & \text{memiliki jari-jari konvergensi } 0.
& &
\end{align*}
\end{example}

\begin{example}
Perhatikan perbedaan antara $\frac{1}{1-z}$ dan deret pangkatnya.  Mari kita
kembangkan $\frac{1}{1-z}$ sebagai deret pangkat di sekitar titik $a \neq 1$.
Misalkan $c \coloneqq \frac{1}{1-a}$, maka
\begin{equation*}
\frac{1}{1-z} =
\frac{c}{1-c(z-a)}
=
c
\sum_{n=0}^\infty c^{n} {(z-a)}^n
=
\sum_{n=0}^\infty \left( \frac{1}{{(1-a)}^{n+1}} \right) {(z-a)}^n .
\end{equation*}
Deret $\sum_{n=0}^\infty c^n {(z-a)}^n$ konvergen jika dan hanya jika
deret di ruas kanan konvergen.  Kita hitung,
\begin{equation*}
\limsup_{n\to\infty}
\sqrt[n]{\sabs{c^n}} = \sabs{c}
= \frac{1}{\sabs{1-a}} .
\end{equation*}
Jari-jari konvergensi deret pangkat tersebut adalah $\sabs{1-a}$, yaitu
jarak dari $1$ ke $a$.  Fungsi $\frac{1}{1-z}$
memiliki representasi deret pangkat
di sekitar setiap $a \neq 1$ sehingga bersifat analitik pada $\C \setminus
\{ 1 \}$.
Domain fungsi tersebut lebih luas daripada daerah
konvergensi deret pangkat yang merepresentasikan fungsi itu di sekitar titik mana pun.
\end{example}

Ternyata,
jika suatu fungsi memiliki representasi deret pangkat yang konvergen ke fungsi tersebut
pada suatu bola,
maka fungsi itu memiliki representasi deret pangkat di sekitar setiap titik dalam bola tersebut.  Kita akan membuktikan
hasil ini nanti.

\subsection{Sifat-sifat fungsi analitik}

\begin{prop}
Jika
\begin{equation*}
f(z) \coloneqq \sum_{n=0}^\infty c_n {(z-a)}^n
\end{equation*}
konvergen pada $B(a,\rho)$ untuk suatu $\rho > 0$, maka
$f \colon B(a,\rho) \to \C$ kontinu.
Khususnya, fungsi analitik bersifat kontinu.
\end{prop}

\begin{proof}
Untuk $z_0 \in B(a,\rho)$, pilih $r < \rho$ sedemikian sehingga $z_0 \in B(a,r)$.
Pada $B(a,r)$,
jumlah parsial (yang kontinu) konvergen secara seragam,
sehingga limitnya, $f|_{B(a,r)}$, kontinu.
Setiap barisan yang konvergen ke
$z_0$ memiliki suatu ekor yang seluruhnya berada dalam bola terbuka $B(a,r)$.
Jadi, $f$ kontinu di $z_0$.
\end{proof}

Dalam \volIref{\corref*{vI-cor:differentiatepowerser} pada Jilid
I}{\corref{cor:differentiatepowerser}}, kita telah membuktikan bahwa kita dapat
mendiferensialkan deret pangkat real suku demi suku.  Dengan kata lain,
kita membuktikan bahwa jika
\begin{equation*}
f(x) \coloneqq \sum_{n=0}^\infty c_n {(x-a)}^n
\end{equation*}
konvergen untuk $x$ real pada suatu interval di sekitar $a \in \R$, maka kita dapat mendiferensialkannya
suku demi suku dan memperoleh deret
\begin{equation*}
f'(x) =
\sum_{n=1}^\infty n c_n {(x-a)}^{n-1}
=
\sum_{n=0}^\infty (n+1)c_{n+1} {(x-a)}^{n} 
\end{equation*}
dengan jari-jari konvergensi yang sama.
Kita hanya membuktikan teorema ini ketika $c_n$ bernilai real.
Namun, untuk $c_n$ kompleks,
kita tulis
$c_n = s_n + i t_n$, dan karena $x$ serta $a$ real,
\begin{equation*}
\sum_{n=0}^\infty c_n {(x-a)}^n
=
\sum_{n=0}^\infty s_n {(x-a)}^n
+
i
\sum_{n=0}^\infty t_n {(x-a)}^n .
\end{equation*}
Kita terapkan teorema tersebut pada bagian real dan
bagian imajiner.

Dengan menerapkan teorema ini berulang kali, kita peroleh bahwa suatu
fungsi analitik diferensiabel tak berhingga kali:
\begin{equation*}
%mbxlatex \begin{aligned}
f^{(\ell)}(x)
%mbxlatex &
=
\sum_{n=\ell}^\infty n(n-1)\cdots(n-\ell+1)c_n {(x-a)}^{n-\ell}
%mbxlatex \\
%mbxlatex &
=
\sum_{n=0}^\infty (n+\ell)(n+\ell-1)\cdots (n+1) c_{n+\ell} {(x-a)}^{n} .
%mbxlatex \end{aligned}
\end{equation*}
Khususnya,
\begin{equation} \label{eq:formulaforpscoeffs}
f^{(\ell)}(a) = \ell! \, c_\ell .
\avoidbreak
\end{equation}
Koefisien-koefisien tersebut ditentukan secara unik oleh turunan
fungsi tersebut, dan sebaliknya.

Di sisi lain, adanya fungsi yang diferensiabel tak berhingga kali
tidak berarti bahwa bilangan $c_n$ yang diperoleh dari
$c_n = \frac{f^{(n)}(0)}{n!}$ menghasilkan deret pangkat yang konvergen.
Terdapat sebuah teorema, yang tidak akan kita buktikan,
yang menyatakan bahwa untuk sebarang barisan $\{ c_n \}_{n=1}^\infty$, terdapat
fungsi $f$ yang diferensiabel tak berhingga kali sedemikian sehingga
$c_n = \frac{f^{(n)}(0)}{n!}$.  Selain itu, sekalipun deret yang diperoleh
konvergen, deret itu belum tentu konvergen ke fungsi semula.
Sebagai contoh,
lihat \volIref{\exerciseref*{vI-exercise:nonanalytic} pada Jilid I}{\exerciseref{exercise:nonanalytic}}:  Fungsi
berikut
\begin{equation*}
f(x) \coloneqq
\begin{cases}
e^{-1/x} & \text{jika } x > 0,\\
0        & \text{jika } x \leq 0,
\end{cases}
\end{equation*}
diferensiabel tak berhingga kali, dan semua turunannya di titik asal bernilai nol.
Jadi,
deretnya di titik asal hanyalah deret nol; walaupun
deret tersebut konvergen, deret itu tidak konvergen ke~$f$ untuk $x > 0$.

\medskip

Kita dapat menerapkan transformasi afin $z \mapsto z+a$ yang
mengubah deret pangkat di sekitar $a$ menjadi deret di sekitar titik asal.
Artinya, jika
\begin{equation*}
f(z) = \sum_{n=0}^\infty c_n {(z-a)}^n,
\qquad \text{kita tinjau} \qquad
f(z+a) = \sum_{n=0}^\infty c_n {z}^n.
\end{equation*}
Karena itu, biasanya
cukup untuk membuktikan hasil mengenai deret pangkat di sekitar titik asal.
Mulai sekarang, demi kesederhanaan, kita sering mengandaikan $a=0$.

\subsection{Deret pangkat sebagai fungsi analitik}

Kita memerlukan teorema tentang pertukaran limit deret, yakni
teorema Fubini untuk penjumlahan.  Untuk deret real, hasil ini merupakan
\volIref{\exerciseref*{vI-exercise:tonellifubiniforsums} pada Jilid
I}{\exerciseref{exercise:tonellifubiniforsums}}, tetapi sekarang kita memiliki
argumen yang lebih elegan.

\begin{thm}[\myindex{Fubini untuk penjumlahan}] \label{thm:fubiniforsums}
Misalkan $\{ a_{k,m} \}_{k=1,m=1}^\infty$ suatu barisan ganda
bilangan kompleks dan andaikan bahwa untuk setiap $k$, deret
\begin{equation*}
\sum_{m=1}^\infty \sabs{a_{k,m}} \qquad \text{konvergen}
\end{equation*}
dan, lebih lanjut, bahwa
\begin{equation*}
\sum_{k=1}^\infty \left( \sum_{m=1}^\infty \sabs{a_{k,m}} \right)
\qquad \text{konvergen}.
\end{equation*}
Maka
\begin{equation*}
\sum_{k=1}^\infty \left( \sum_{m=1}^\infty a_{k,m} \right)
=
\sum_{m=1}^\infty \left( \sum_{k=1}^\infty a_{k,m} \right) ,
\end{equation*}
dan semua deret yang terlibat konvergen.
\end{thm}

\begin{proof}
Misalkan $E$ himpunan $\{ \nicefrac{1}{n} : n \in \N \} \cup \{ 0 \}$,
dan pandang himpunan tersebut sebagai ruang metrik dengan metrik yang diwarisi dari $\R$.
Definisikan barisan fungsi $f_k \colon E \to \C$
dengan
\begin{equation*}
f_k(\nicefrac{1}{n}) \coloneqq \sum_{m=1}^n a_{k,m}
\qquad
\text{dan}
\qquad
f_k(0) \coloneqq \sum_{m=1}^\infty a_{k,m} .
\end{equation*}
Karena deret tersebut konvergen, setiap $f_k$ kontinu di $0$
(karena 0 adalah satu-satunya titik akumulasi, fungsi-fungsi tersebut kontinu di setiap titik
dalam $E$, tetapi kita tidak memerlukan fakta itu).
Untuk setiap $x \in E$, kita memiliki
\begin{equation*}
\babs{f_k(x)} \leq \sum_{m=1}^\infty \sabs{a_{k,m}} .
\end{equation*}
Karena $\sum_k \sum_m \sabs{a_{k,m}}$ konvergen (dan tidak bergantung pada
$x$), kita tahu bahwa
\begin{equation*}
\sum_{k=1}^n f_k(x)
\end{equation*}
konvergen secara seragam pada $E$.  Definisikan
\begin{equation*}
g(x) \coloneqq \sum_{k=1}^\infty f_k(x) ,
\end{equation*}
yang karena itu merupakan fungsi kontinu di $0$.
Dengan demikian,
\begin{equation*}
\begin{split}
\sum_{k=1}^\infty \left( \sum_{m=1}^\infty a_{k,m} \right)
& =
\sum_{k=1}^\infty f_k(0)
= g(0)
= \lim_{n\to\infty} g(\nicefrac{1}{n}) \\
&= 
\lim_{n\to\infty}\sum_{k=1}^\infty f_k(\nicefrac{1}{n})
= 
\lim_{n\to\infty}\sum_{k=1}^\infty \sum_{m=1}^n a_{k,m} \\
&= 
\lim_{n\to\infty}\sum_{m=1}^n \sum_{k=1}^\infty a_{k,m}
= 
\sum_{m=1}^\infty \left( \sum_{k=1}^\infty a_{k,m} \right) . \qedhere
\end{split}
\end{equation*}
\end{proof}

Sekarang kita buktikan bahwa jika suatu deret konvergen ke sebuah fungsi
pada suatu interval, kita dapat mengembangkan fungsi tersebut di sekitar setiap titik.

\begin{thm}[Teorema Taylor untuk fungsi analitik real]
\index{Teorema Taylor!analitik real}%
\label{thm:tayloranal}
Misalkan
\begin{equation*}
f(x) \coloneqq \sum_{k=0}^\infty a_k x^k
\end{equation*}
merupakan suatu deret pangkat yang konvergen pada $(-\rho,\rho)$ untuk suatu $\rho > 0$.  Untuk setiap $a \in
(-\rho,\rho)$
dan setiap $x$ yang memenuhi $\sabs{x-a} < \rho-\sabs{a}$, berlaku
\begin{equation*}
f(x) =
\sum_{k=0}^\infty \frac{f^{(k)}(a)}{k!} {(x-a)}^{k} .
\end{equation*}
\end{thm}

Deret pangkat yang berpusat di $a$ tentu saja dapat konvergen pada interval yang lebih besar, tetapi
konvergensi pada interval di atas dijamin.  Interval tersebut adalah interval simetris terbesar yang berpusat di
$a$ dan termuat dalam $(-\rho,\rho)$.

\begin{proof}
Untuk $a$ dan $x$ sebagaimana dalam teorema,
tuliskan
\begin{equation*}
\begin{split}
f(x) &= \sum_{k=0}^\infty a_k {\bigl((x-a)+a\bigr)}^k \\
&= \sum_{k=0}^\infty a_k \sum_{m=0}^k \binom{k}{m} a^{k-m} {(x-a)}^m .
\end{split}
\end{equation*}
Definisikan $c_{k,m} \coloneqq a_k \binom{k}{m} a^{k-m}$ jika $m \leq k$ dan $0$ jika $m >
k$.  Maka
\begin{equation} \label{eq:tsproof}
f(x) = \sum_{k=0}^\infty \, \sum_{m=0}^\infty c_{k,m} {(x-a)}^m .
\end{equation}
Mari kita tunjukkan bahwa jumlah ganda tersebut konvergen mutlak.
\begin{equation*}
\begin{split}
\sum_{k=0}^\infty \, \sum_{m=0}^\infty \babs{ c_{k,m} {(x-a)}^m}
& = \sum_{k=0}^\infty \, \sum_{m=0}^k \abs{ a_k \binom{k}{m} a^{k-m} {(x-a)}^m }
\\
& = \sum_{k=0}^\infty \sabs{a_k} \sum_{m=0}^k \binom{k}{m} \sabs{a}^{k-m}
{\sabs{x-a}}^m  \\
& = \sum_{k=0}^\infty \sabs{a_k} {\bigl(\sabs{x-a}+\sabs{a}\bigr)}^k ,
\end{split}
\end{equation*}
dan deret ini konvergen selama
$(\sabs{x-a}+\sabs{a}) < \rho$, atau dengan kata lain jika
$\sabs{x-a} < \rho-\sabs{a}$.

Dengan menggunakan \thmref{thm:fubiniforsums},
tukar urutan penjumlahan dalam \eqref{eq:tsproof}; maka
deret berikut konvergen ketika $\sabs{x-a} < \rho-\sabs{a}$:
\begin{equation*}
f(x) =
\sum_{k=0}^\infty \, \sum_{m=0}^\infty c_{k,m} {(x-a)}^m
=
\sum_{m=0}^\infty
\left( \sum_{k=0}^\infty
c_{k,m} \right) {(x-a)}^m .
\end{equation*}
Rumus yang dinyatakan melalui turunan di $a$ diperoleh dengan
mendiferensialkan deret tersebut untuk memperoleh \eqref{eq:formulaforpscoeffs}.
\end{proof}

Perhatikan bahwa jika suatu deret konvergen untuk peubah real $x \in (a-\rho,a+\rho)$, deret itu juga konvergen
untuk semua bilangan kompleks dalam $B(a,\rho)$.
Kita memperoleh korolari berikut, yang menyatakan bahwa fungsi yang didefinisikan
oleh deret pangkat bersifat analitik.

\begin{cor} \label{cor:powerseranalytic}
Untuk setiap $a \in \C$,
jika $\sum_{k=0}^\infty c_k {(z-a)}^k$ konvergen ke $f(z)$ dalam $B(a,\rho)$ dan $b \in
B(a,\rho)$,
maka terdapat deret pangkat
$\sum_{k=0}^\infty d_k {(z-b)}^k$ yang konvergen ke $f(z)$ dalam $B(b,\rho-\sabs{b-a})$.
\end{cor}

\begin{proof}
Tanpa mengurangi keumuman, andaikan $a=0$ dan $b \neq 0$.  Kita dapat melakukan rotasi sehingga $b$ bernilai real, tetapi
karena hal itu lebih sulit dibayangkan, mari kita lakukan secara eksplisit.
Misalkan $\alpha \coloneqq \frac{\bar{b}}{\sabs{b}}$.
Perhatikan bahwa
\begin{equation*}
\abs{\nicefrac{1}{\alpha}} = \sabs{\alpha}
= 1 .
\end{equation*}
Karena itu, deret
$\sum_{k=0}^\infty c_k {(\nicefrac{z}{\alpha})}^k = 
\sum_{k=0}^\infty c_k \alpha^{-k} {z}^k$
konvergen ke $f(\nicefrac{z}{\alpha})$ dalam $B(0,\rho)$.
Ketika $z=x$ bernilai real,
kita menerapkan \thmref{thm:tayloranal} pada titik $\sabs{b}$ dan memperoleh
sebuah deret yang konvergen
ke $f(\nicefrac{z}{\alpha})$ pada $B(\sabs{b},\rho-\sabs{b})$.
Artinya, terdapat deret konvergen
\begin{equation*}
f(\nicefrac{z}{\alpha}) =
\sum_{k=0}^\infty a_k {\bigl(z - \sabs{b}\bigr)}^k .
\end{equation*}
Dengan menggunakan $\alpha b = \sabs{b}$, kita peroleh
\begin{equation*}
f(z) = f(\nicefrac{\alpha z}{\alpha}) =
\sum_{k=0}^\infty a_k {(\alpha z - \sabs{b})}^k 
=
\sum_{k=0}^\infty a_k\alpha^k {\bigl(z - \nicefrac{\sabs{b}}{\alpha}\bigr)}^k
=
\sum_{k=0}^\infty a_k\alpha^k {(z - b)}^k ,
\end{equation*}
dan deret ini konvergen untuk semua $z$ sedemikian sehingga
$\bigl\lvert \alpha z-\sabs{b}\bigr\rvert < \rho-\sabs{b}$
atau $\sabs{z - b} < \rho-\sabs{b}$.
\end{proof}

Di atas telah kita buktikan bahwa deret pangkat yang konvergen mendefinisikan
fungsi analitik pada daerah konvergensinya.  Sebelumnya juga telah kita tunjukkan bahwa
$\frac{1}{1-z}$ bersifat analitik di setiap titik selain $z=1$.

Perhatikan bahwa meskipun suatu fungsi analitik real bersifat analitik pada
seluruh garis real, hal itu belum tentu berarti fungsi tersebut memiliki representasi deret pangkat
yang konvergen di semua titik.  Sebagai contoh, fungsi
\begin{equation*}
f(x) = \frac{1}{1+x^2}
\end{equation*}
merupakan fungsi analitik real pada $\R$ (latihan).  Deret
pangkat di sekitar titik asal yang konvergen ke $f$
memiliki jari-jari konvergensi tepat $1$.
Dapatkah Anda melihat alasannya? (latihan)

\subsection{Teorema identitas untuk fungsi analitik}

\begin{lemma}
Misalkan $f(z) = \sum_{k=0}^\infty a_k z^k$ adalah deret pangkat konvergen dan
$\{ z_n \}_{n=1}^\infty$ adalah barisan bilangan kompleks tak nol yang konvergen ke 0,
sedemikian sehingga $f(z_n) = 0$ untuk semua $n$.  Maka $a_k = 0$ untuk setiap~$k$.
\end{lemma}

\begin{proof}
Berdasarkan kekontinuan, kita tahu bahwa $f(0) = 0$, sehingga $a_0 = 0$.
Misalkan terdapat suatu $a_k$ yang tak nol.
Ambil $m$ sebagai indeks terkecil sedemikian sehingga $a_m \neq 0$.  Maka
\begin{equation*}
f(z) = \sum_{k=m}^\infty a_k z^k = 
z^m \sum_{k=m}^\infty a_k z^{k-m} =
z^m \sum_{k=0}^\infty a_{k+m} z^{k} .
\end{equation*}
Tuliskan $g(z) = \sum_{k=0}^\infty a_{k+m} z^{k}$ (deret ini konvergen
pada himpunan yang sama seperti deret untuk $f$).  Fungsi $g$ kontinu dan
$g(0) = a_m \neq 0$.  Jadi, terdapat suatu $\delta > 0$ sedemikian sehingga
$g(z) \neq 0$ untuk semua $z \in B(0,\delta)$.  Karena $f(z) = z^m g(z)$,
satu-satunya titik di $B(0,\delta)$ tempat $f(z) = 0$ adalah $z=0$, tetapi hal
ini bertentangan dengan asumsi bahwa $f(z_n) = 0$ untuk semua $n$.
\end{proof}

Ingatlah bahwa dalam ruang metrik $X$, sebuah \emph{titik akumulasi}
(atau terkadang \emph{titik limit}) dari suatu himpunan
$E$ adalah titik $p \in X$ sedemikian sehingga
$B(p,\epsilon) \setminus \{ p \}$ memuat titik-titik dari $E$
untuk setiap $\epsilon > 0$.

\begin{thm}[Teorema identitas]
\label{thm:identityanalytic}%
\index{teorema identitas}%
Misalkan $U \subset \C$ terbuka dan terhubung.  Jika $f \colon U \to \C$
dan $g \colon U \to \C$ merupakan fungsi-fungsi analitik yang bernilai
sama pada suatu himpunan $E \subset U$, dan $E$ memiliki titik akumulasi
di $U$, maka $f(z) = g(z)$ untuk semua $z \in U$.
\end{thm}

Dalam penerapan yang paling umum dari teorema ini, $E$ adalah himpunan terbuka
atau mungkin sebuah kurva.

\begin{proof}
Dengan mengganti $E$, jika perlu, dengan himpunan semua titik $z \in U$ yang
memenuhi $g(z)=f(z)$, kita boleh menganggap $E$ adalah himpunan tersebut.
Perhatikan bahwa $E$ harus tertutup relatif terhadap $U$ karena $f$ dan $g$
kontinu.

Misalkan $p \in U$ adalah suatu titik akumulasi dari $E$.  Dengan mengganti
peubah $w=z-p$ serta mentranslasikan domain dan kedua fungsi, kita boleh
mengambil $p=0$.  Di sekitar $0$, kita memiliki ekspansi
\begin{equation*}
f(z) = \sum_{k=0}^\infty a_k {z}^k 
\qquad
\text{dan}
\qquad
g(z) = \sum_{k=0}^\infty b_k {z}^k ,
\end{equation*}
yang konvergen dalam suatu bola $B(0,\rho)$.  Oleh karena itu, deret
\begin{equation*}
0 = f(z)-g(z) = 
\sum_{k=0}^\infty (a_k-b_k) z^k
\end{equation*}
konvergen dalam $B(0,\rho)$.  Karena $0$ adalah titik akumulasi dari $E$,
terdapat barisan titik tak nol $\{ z_n \}_{n=1}^\infty$ yang konvergen ke $0$
sedemikian sehingga $f(z_n) -g(z_n) = 0$ untuk semua $n$.  Jadi, berdasarkan lema di atas,
$a_k = b_k$ untuk semua $k$.  Oleh karena itu, $B(0,\rho) \subset E$.

Dengan demikian, himpunan semua titik akumulasi dari $E$ terbuka relatif
terhadap $U$.  Himpunan itu juga tertutup relatif terhadap $U$: limit di $U$
dari titik-titik akumulasi $E$ berada di $E$ karena $E$ tertutup relatif
terhadap $U$, dan jelas merupakan titik akumulasi dari $E$.
Karena $U$ terhubung, himpunan semua titik akumulasi dari $E$
sama dengan $U$, atau dengan kata lain $E = U$.
\end{proof}

Dengan membatasi perhatian kita pada $x$ real, kita memperoleh teorema yang sama
untuk himpunan bagian terbuka dan terhubung dari $\R$, yaitu interval terbuka.

\subsection{Latihan}

\begin{exercise}
Misalkan
\begin{equation*}
a_{k,m} \coloneqq
\begin{cases}
1        & \text{jika } k=m,\\
-2^{k-m} & \text{jika } k<m,\\
0        & \text{jika } k>m.
\end{cases}
\end{equation*}
Hitung (atau tunjukkan bahwa limitnya tidak ada):
\\
a)~$\displaystyle \sum_{m=1}^\infty \sabs{a_{k,m}}~$ untuk setiap $k$,
%mbxSTARTIGNORE
\hspace{\fill}
%mbxENDIGNORE
%mbxlatex \qquad
b)~$\displaystyle \sum_{k=1}^\infty \sabs{a_{k,m}}~$ untuk setiap $m$,
%mbxSTARTIGNORE
\hspace{\fill}
%mbxENDIGNORE
%mbxlatex \qquad
c)~$\displaystyle \sum_{k=1}^\infty \sum_{m=1}^\infty \sabs{a_{k,m}}$,
%mbxSTARTIGNORE
\hspace{\fill}
%mbxENDIGNORE
%mbx <rabr/>
d)~$\displaystyle \sum_{k=1}^\infty \sum_{m=1}^\infty a_{k,m}$,
%mbxSTARTIGNORE
\hspace{\fill}
%mbxENDIGNORE
%mbxlatex \qquad
e)~$\displaystyle \sum_{m=1}^\infty \sum_{k=1}^\infty a_{k,m}$.
\\
Petunjuk: Teorema Fubini untuk deret tidak berlaku.
Bahkan, jawaban untuk d) dan e) berbeda.
\end{exercise}

\begin{exercise}
\pagebreak[2]
Misalkan $f(x) \coloneqq \frac{1}{1+x^2}$.  Buktikan bahwa
\begin{enumerate}[a)]
\item
$f$ merupakan fungsi analitik pada seluruh $\R$
dengan menemukan deret pangkat untuk $f$ di setiap $a \in \R$,
\item
jari-jari konvergensi dari
deret pangkat untuk $f$ di titik asal adalah 1.
\end{enumerate}
\end{exercise}


\begin{exercise}
\pagebreak[2]
Misalkan $f \colon \C \to \C$ analitik dan bukan fungsi nol.  Tunjukkan bahwa untuk setiap
$n$, hanya terdapat berhingga banyak titik nol dari $f$ di $B(0,n)$, yaitu,
$f^{-1}(0) \cap B(0,n)$ berhingga untuk setiap $n$.
\end{exercise}

\begin{exercise}
Misalkan $U \subset \C$ terbuka dan terhubung, $0 \in U$, dan $f \colon U
\to \C$ analitik.  Dengan memandang $f$ sebagai fungsi dari peubah real $x$
di titik asal, misalkan $f^{(n)}(0) = 0$ untuk semua $n$.  Tunjukkan bahwa $f(z) = 0$
untuk semua $z \in U$.
\end{exercise}

\begin{exercise}
Misalkan $U \subset \C$ terbuka dan terhubung, $0 \in U$, dan $f \colon U
\to \C$ analitik.  Untuk $x$ dan $y$ real yang cukup dekat dengan $0$ sehingga
$x \in U$ dan $iy \in U$, misalkan $h(x) \coloneqq f(x)$ dan
$g(y) \coloneqq -i \, f(iy)$.
Tunjukkan bahwa $h$ dan $g$ diferensiabel tak berhingga kali di titik asal dan
$h'(0) = g'(0)$.
\end{exercise}

\begin{exercise}
Misalkan suatu fungsi $f$ analitik pada suatu lingkungan dari titik asal,
dan terdapat suatu $M$
sedemikian sehingga $\sabs{f^{(n)}(0)} \leq M$ untuk semua $n$.
Buktikan bahwa deret pangkat $f$ di titik asal konvergen untuk semua $z \in \C$.
\end{exercise}

\begin{exercise}
Misalkan $f(z) \coloneqq \sum_{n=0}^\infty c_n z^n$ dengan jari-jari konvergensi 1.  Misalkan $f(0)
= 0$, tetapi $f$ bukan fungsi nol.
Tunjukkan bahwa terdapat suatu $k \in \N$ dan suatu
deret pangkat konvergen $g(z) \coloneqq \sum_{n=0}^\infty d_n z^n$ dengan jari-jari konvergensi 1
sedemikian sehingga $f(z) = z^k g(z)$ untuk semua $z \in B(0,1)$, dan $g(0) \neq 0$.
\end{exercise}

\begin{exercise}
Misalkan $U \subset \C$ terbuka dan terhubung.  Misalkan
$f \colon U \to \C$ analitik, $U \cap \R \neq \emptyset$, dan
$f(x) = 0$ untuk semua $x \in U \cap \R$.  Tunjukkan bahwa $f(z) = 0$ untuk semua $z \in
U$.
\end{exercise}

\begin{exercise}
Untuk $\alpha \in \C$ dan $k=0,1,2,3\ldots$, definisikan
\begin{equation*}
\binom{\alpha}{k} \coloneqq \frac{\alpha(\alpha-1)\cdots(\alpha-k+1)}{k!} .
\end{equation*}
\begin{enumerate}[a)]
\item
Tunjukkan bahwa deret
\begin{equation*}
f(z) \coloneqq \sum_{k=0}^\infty \binom{\alpha}{k} z^k
\end{equation*}
konvergen setiap kali $\sabs{z} < 1$.
Bahkan, buktikan bahwa untuk $\alpha = 0,1,2,3,\ldots$ jari-jari
konvergensinya adalah $\infty$, dan untuk semua $\alpha$ lainnya
jari-jari konvergensinya adalah 1.
\item
Tunjukkan bahwa untuk $x \in \R$, $\sabs{x} < 1$, kita mempunyai
\begin{equation*}
(1+x) f'(x) = \alpha f(x) ,
\end{equation*}
yang berarti bahwa $f(x) = (1+x)^\alpha$.
\end{enumerate}
\end{exercise}

\begin{exercise}
Misalkan $f \colon \C \to \C$ analitik dan misalkan untuk suatu
interval terbuka $(a,b) \subset \R$, $f$ bernilai real pada $(a,b)$.  Tunjukkan
bahwa $f$ bernilai real pada $\R$.
\end{exercise}

\begin{exercise}
Misalkan $\D \coloneqq B(0,1)$ adalah cakram satuan.  Misalkan
$f \colon \D \to \C$ analitik dengan deret pangkat
$\sum_{n=0}^\infty c_n z^n$.  Misalkan $\sabs{c_n} \leq 1$ untuk semua $n$.  Buktikan bahwa untuk
semua $z \in \D$, kita mempunyai
$\babs{f(z)} \leq \frac{1}{1-\sabs{z}}$.
\end{exercise}

%%%%%%%%%%%%%%%%%%%%%%%%%%%%%%%%%%%%%%%%%%%%%%%%%%%%%%%%%%%%%%%%%%%%%%%%%%%%%%

\sectionnewpage
\section{Eksponensial kompleks dan fungsi trigonometri}
\label{sec:complexexp}

%mbxINTROSUBSECTION

\sectionnotes{1 kuliah}

\subsection{Eksponensial kompleks}

Definisikan
\begin{equation*}
E(z) \coloneqq \sum_{k=0}^\infty \frac{1}{k!} z^k .
\end{equation*}
Deret ini konvergen untuk setiap $z \in \C$.
Jadi, menurut
\corref{cor:powerseranalytic}, $E$ analitik pada $\C$. Kita perhatikan bahwa $E(0) = 1$
dan, untuk $z=x \in \R$, berlaku $E(x) \in \R$. Dengan tetap mengambil $x$ real,
perhitungan langsung menunjukkan
\begin{equation*}
\frac{d}{dx} \bigl( E(x) \bigr) = E(x) .
\end{equation*}
Dalam \volIref{\sectionref*{vI-sec:logandexp} pada volume I}{\sectionref{sec:logandexp}} (atau dengan Teorema Picard), kita membuktikan bahwa
satu-satunya fungsi pada garis real yang memenuhi $E' = E$ dan
$E(0) = 1$ adalah fungsi eksponensial. Dengan kata lain, untuk $x \in \R$, $e^x = E(x)$.

Untuk bilangan kompleks $z$, kita definisikan\glsadd{not:complexexp}
\begin{equation*}
e^z \coloneqq E(z) = 
\sum_{k=0}^\infty \frac{1}{k!} z^k .
\end{equation*}
Pada garis real, definisi baru ini sama dengan definisi kita sebelumnya.
Lihat \figureref{fig:complexexpgraphs}. Perhatikan bahwa pada arah $x$
(arah real),
grafik berperilaku seperti eksponensial real, sedangkan pada arah $y$
(arah imajiner) grafik berosilasi.

\begin{myfigureht}
\myincludepdft{real_imag_exp}{%
Grafik permukaan.
Untuk x tetap, kedua grafik berbentuk sinusoidal pada arah y,
dengan amplitudo osilasi yang semakin besar ketika x semakin besar.
Untuk x negatif, grafik bahkan begitu kecil pada skala ini sehingga
osilasinya tidak dapat terlihat.
Pada grafik kanan, gelombang sinusoidal digeser. Pada grafik kiri,
irisan tepat di atas sumbu x naik secara eksponensial, sedangkan pada grafik
kanan irisan tepat di atas sumbu x selalu bernilai nol; dengan kata lain,
sumbu x termuat dalam grafik kanan. Irisan grafik di atas sumbu x ditandai
dengan garis tebal.}
\caption{Grafik bagian real (kiri) dan bagian imajiner (kanan)
dari eksponensial kompleks $e^z = e^{x+iy}$. Sumbu $x$ membentang dari $-4$ hingga
$4$, sumbu $y$ membentang dari $-6$ hingga $6$, dan sumbu vertikal membentang dari
$-e^{4} \approx -54.6$
hingga
$e^{4} \approx 54.6$. Grafik eksponensial real ($y=0$)
ditandai dengan garis tebal.\label{fig:complexexpgraphs}}
\end{myfigureht}

\begin{prop}[Hukum eksponen\index{hukum eksponen}]
Misalkan $z,w \in \C$ adalah bilangan kompleks. Maka
\begin{equation*}
e^{z+w} = e^z e^w.
\end{equation*}
\end{prop}

\begin{proof}
Kita sudah mengetahui bahwa kesamaan
$e^{x+y} = e^x e^y$ berlaku untuk semua
bilangan real $x$ dan $y$.
Untuk setiap $y \in \R$ yang tetap, pandang kedua ruas sebagai
fungsi dari $x$ dan terapkan teorema identitas
(\thmref{thm:identityanalytic}) untuk memperoleh
$e^{z+y} = e^ze^y$ bagi setiap $z \in \C$. Dengan menetapkan sembarang $z \in \C$,
kita memperoleh
$e^{z+y} = e^ze^y$ bagi setiap $y \in \R$. Sekali lagi menurut teorema identitas,
$e^{z+w} = e^z e^w$
bagi setiap $w \in \C$.
\end{proof}

Konsekuensi langsung dari proposisi ini adalah $e^z \neq 0$ untuk setiap $z \in \C$,
karena $e^z e^{-z} = e^{z-z} = 1$.
Perhitungan ini berarti bahwa ${(e^z)}^{-1} = e^{-z}$.
Dengan menggabungkan fakta tersebut dan hukum eksponen, kita memperoleh
\begin{equation*}
{(e^z)}^{n} = e^{nz} \qquad \text{untuk setiap } n \in \Z.
\end{equation*}
Konsekuensi lain yang sedikit lebih rumit adalah bahwa kita
dapat menghitung deret pangkat untuk fungsi eksponensial di sembarang titik $a \in
\C$:
\begin{equation*}
e^z = e^a e^{z-a} = \sum_{k=0}^\infty \frac{e^a}{k!} {(z-a)}^k .
\end{equation*}

\subsection{Fungsi trigonometri dan $\pi$}

Kini akhirnya kita dapat mendefinisikan \emph{\myindex{fungsi sinus}} dan \emph{\myindex{fungsi kosinus}}
melalui persamaan
\begin{equation*}
e^{x+iy} = e^x \bigl( \cos(y) + i \sin(y) \bigr) .
\end{equation*}
Lebih umum, kita mendefinisikan fungsi sinus dan fungsi kosinus untuk setiap bilangan kompleks $z$:
\glsadd{not:sin}\glsadd{not:cos}
\begin{equation*}
\cos(z) \coloneqq \frac{e^{iz} + e^{-iz}}{2}
\qquad\text{dan}\qquad
\sin(z) \coloneqq \frac{e^{iz} - e^{-iz}}{2i} .
\end{equation*}

Mari kita gunakan definisi ini untuk membuktikan sifat-sifat umum
fungsi sinus dan fungsi kosinus. Dalam prosesnya, kita juga mendefinisikan
bilangan $\pi$.

\begin{prop}
Fungsi sinus dan fungsi kosinus memiliki sifat-sifat berikut:
\begin{enumerate}[(i)]
\item Untuk setiap $z \in \C$,\index{rumus Euler}
\begin{equation*}
e^{iz} = \cos(z) + i\sin(z) \qquad
\text{(rumus Euler)}.
\end{equation*}
\item $\cos(0) = 1$, $\sin(0) = 0$.
\item Untuk setiap $z \in \C$,
\begin{equation*}
\cos(-z) = \cos(z), \qquad
\sin(-z) = -\sin(z).
\end{equation*}
\item Untuk setiap $z \in \C$,
\begin{equation*}
\cos(z) = \sum_{k=0}^\infty \frac{{(-1)}^k}{(2k)!} z^{2k} ,
\qquad
\sin(z) = \sum_{k=0}^\infty \frac{{(-1)}^k}{(2k+1)!} z^{2k+1} .
\end{equation*}
\item Untuk setiap $x \in \R$
\begin{equation*}
\cos(x) = \Re (e^{ix})
\qquad\text{dan}\qquad
\sin(x) = \Im (e^{ix}) .
\end{equation*}
\item Untuk setiap $z \in \C$,
\begin{equation*}
{\bigl( \cos(z) \bigr)}^2 + {\bigl( \sin(z) \bigr)}^2 = 1 .
\end{equation*}
\item Untuk setiap $x \in \R$,
\begin{equation*}
\babs{\sin(x)} \leq 1, \qquad \babs{\cos(x)} \leq 1 .
\end{equation*}
\item Untuk setiap $x \in \R$,
\begin{equation*}
\frac{d}{dx} \bigl[ \cos(x) \bigr] = -\sin(x)
\qquad \text{dan} \qquad
\frac{d}{dx} \bigl[ \sin(x) \bigr] = \cos(x) .
\end{equation*}
\item Untuk setiap $x \geq 0$,
\begin{equation*}
\sin(x) \leq x .
\end{equation*}
\item
Terdapat $x > 0$ sedemikian sehingga $\cos(x) = 0$. Kita definisikan
\glsadd{not:pi}
\begin{equation*}
\pi \coloneqq 2 \, \inf \{ x > 0 : \cos(x) = 0 \} .
\end{equation*}
\item
Untuk setiap $z \in \C$,
\begin{equation*}
e^{2\pi i} = 1 \qquad \text{dan} \qquad e^{z + i 2\pi} = e^z.
\end{equation*}
\item
Fungsi sinus dan fungsi kosinus periodik dengan periode $2\pi$ dan tidak memiliki periode positif
yang lebih kecil. Artinya, $2\pi$ adalah bilangan positif terkecil sedemikian sehingga untuk setiap $z \in \C$,
\begin{equation*}
\sin(z+2\pi) = \sin(z)
\qquad \text{dan} \qquad
\cos(z+2\pi) = \cos(z) .
\end{equation*}
\item
Fungsi $x \mapsto e^{ix}$ merupakan pemetaan bijektif dari $[0,2\pi)$
ke himpunan semua $z \in \C$ sedemikian sehingga $\sabs{z} = 1$.
\end{enumerate}
\end{prop}

Proposisi ini langsung menyiratkan bahwa $\sin(x)$ dan $\cos(x)$ bernilai
real apabila $x$ real.

\begin{proof}
Tiga butir pertama langsung mengikuti definisi.
Perhitungan deret pangkat untuk kedua fungsi tersebut diserahkan sebagai latihan.
Karena pemetaan konjugasi kompleks bersifat kontinu, definisi
$e^z$ menyiratkan
$\overline{e^z} = e^{\bar{z}}$. Jika
$x$ real,
\begin{equation*}
\overline{e^{ix}} = e^{-ix} .
\end{equation*}
Jadi, untuk $x$ real,
$\cos(x) =
\frac{e^{ix}+e^{-ix}}{2} =
\frac{e^{ix}+\overline{e^{ix}}}{2} =
\Re (e^{ix})$
dan, secara serupa, $\sin(x) = \Im (e^{ix})$.

Untuk $x$ real, kita hitung
\begin{equation*}
1 =  e^{ix} e^{-ix}
= e^{ix} \, \overline{e^{ix}}
= \sabs{e^{ix}}^2
= \babs{\cos(x) + i \sin(x)}^2
= {\bigl( \cos(x) \bigr)}^2 + {\bigl( \sin(x) \bigr)}^2 .
\end{equation*}
Perhitungan yang sedikit lebih rumit menunjukkan fakta ini untuk
bilangan kompleks; lihat \exerciseref{exercise:cossinidentity}.
Khususnya, $e^{ix}$ bersifat \emph{\myindex{bermodulus satu}} untuk $x$ real;
nilai-nilainya terletak pada lingkaran satuan.
Kuadrat suatu bilangan real selalu nonnegatif:
\begin{equation*}
{\bigl(\sin(x)\bigr)}^2 = 1-{\bigl(\cos(x)\bigr)}^2 \leq 1 .
\end{equation*}
Jadi, $\babs{\sin(x)} \leq 1$ dan, secara serupa,
$\babs{\cos(x)} \leq 1$.

Perhitungan kedua turunan tersebut kita serahkan kepada pembaca sebagai latihan.
Mari kita buktikan bahwa $\sin(x) \leq x$ untuk $x \geq 0$.
Pandang
$f(x) \coloneqq x-\sin(x)$ dan turunkan:
\begin{equation*}
f'(x) = \frac{d}{dx} \bigl[ x - \sin(x) \bigr]
=
1 -\cos(x) \geq 0 ,
\end{equation*}
untuk setiap $x \in \R$ karena $\babs{\cos(x)} \leq 1$.
Dengan kata lain, $f$ monoton naik dan $f(0) = 0$.
Jadi, $f$ nonnegatif ketika $x \geq 0$; oleh karena itu, $\sin(x) \leq x$.

Selanjutnya, kita klaim bahwa terdapat $x$ positif sedemikian sehingga $\cos(x) = 0$.
Karena $\cos(0) = 1 > 0$, berlaku $\cos(x) > 0$
untuk $x$ yang cukup dekat dengan $0$. Definisikan
$S \coloneqq \{ t > 0 : \cos(x) > 0 \text{ untuk setiap } x \in [0,t) \}$.
Himpunan $S$ tak kosong. Pilih $y_0 \in S$ dan $a \in (0,y_0)$.
Karena $\cos(x) > 0$ pada $[0,y_0)$, fungsi $\sin(x)$ naik tegas
pada $[0,y_0)$. Karena $\sin(0) = 0$, kita memperoleh $\sin(a) > 0$.
Sekarang, ambil sembarang $y \in S$ dengan $y > a$. Pada $[0,y)$,
fungsi $\sin$ naik tegas. Menurut teorema nilai rata-rata,
terdapat $c \in (a,y)$ sedemikian sehingga
\begin{equation*}
2 \geq \cos(a)-\cos(y) = \sin(c)(y-a) \geq \sin(a)(y-a) .
\end{equation*}
Karena $\sin(a) > 0$, kita memperoleh
\begin{equation*}
y \leq \frac{2}{\sin(a)} + a .
\end{equation*}
Batas ini berlaku bagi setiap $y \in S$ dengan $y > a$, sedangkan anggota $S$
yang memenuhi $y \leq a$ sudah terbatas. Jadi, $S$ terbatas di atas. Definisikan
$y \coloneqq \sup S$. Dari sifat supremum, $\cos(x) > 0$ untuk setiap $x \in [0,y)$,
sehingga kekontinuan memberikan $\cos(y) \geq 0$. Jika $\cos(y) > 0$, kekontinuan
akan memberikan suatu $t > y$ dengan $t \in S$, bertentangan dengan definisi supremum.
Oleh karena itu, $\cos(y) = 0$. Karena $\cos(x) > 0$ untuk setiap $x \in [0,y)$,
$y$ adalah nilai positif terkecil sedemikian sehingga $\cos(y) = 0$. Seperti telah disebutkan,
$\pi$ didefinisikan sebagai $2y$.

Karena $\cos(\nicefrac{\pi}{2}) = 0$, kita peroleh
${\bigl(\sin(\nicefrac{\pi}{2})\bigr)}^2 = 1$.
Karena $\sin$ positif pada $(0,\nicefrac{\pi}{2})$, kita mempunyai
$\sin(\nicefrac{\pi}{2}) = 1$.
Dengan demikian,
\begin{equation*}
e^{i \pi /2} = i ,
\end{equation*}
dan, menurut hukum eksponen,
\begin{equation*}
e^{i \pi} = -1 ,
\qquad 
e^{i 2\pi} = 1 .
\end{equation*}
Jadi, $e^{i2\pi} = 1 = e^0$. Hukum eksponen juga menyatakan
\begin{equation*}
e^{z+i2\pi} = e^z e^{i2\pi} = e^z
\end{equation*}
untuk setiap $z \in \C$. Kita juga langsung memperoleh
$\cos(z+2\pi) = \cos(z)$ dan $\sin(z+2\pi) = \sin(z)$.
Jadi, $\sin$ dan $\cos$ periodik dengan periode $2\pi$.

Kita klaim bahwa $\sin$ dan $\cos$ tidak periodik dengan periode positif yang lebih kecil.
Untuk itu, ambil sembarang $x$ yang memenuhi $0 < x \leq 2\pi$ dan
$e^{ix} = 1$. Kita akan menunjukkan bahwa $x = 2\pi$.
Menurut hukum eksponen,
\begin{equation*}
{\bigl(e^{ix/4}\bigr)}^4 = 1 .
\end{equation*}
Jika $e^{ix/4} = a+ib$, maka
\begin{equation*}
{(a+ib)}^4
=a^4-6a^2b^2+b^4 + i\bigl(4ab(a^2-b^2)\bigr)
=1 .
\end{equation*}
Karena $0 < \nicefrac{x}{4} \leq \nicefrac{\pi}{2}$, kita mempunyai
$b = \sin(\nicefrac{x}{4}) > 0$.
Selain itu, $a = \cos(\nicefrac{x}{4}) \geq 0$. Bagian imajiner pada persamaan di atas
harus nol. Karena $b > 0$, maka $a = 0$ atau $a^2 = b^2$.
Jika $a^2=b^2$, maka
$a^4-6a^2b^2+b^4 = -4a^4 < 0$, dan khususnya tidak sama dengan 1.
Oleh karena itu, $a=0$; dalam hal ini $\nicefrac{x}{4} = \nicefrac{\pi}{2}$.
Jadi, $x = 2\pi$.

Sekarang, misalkan $T > 0$ merupakan periode fungsi sinus. Maka, untuk setiap $x \in \R$,
$\sin(x+T)=\sin(x)$. Dengan menurunkan identitas ini terhadap $x$, kita memperoleh
$\cos(x+T)=\cos(x)$. Jika $T$ merupakan periode fungsi kosinus, penalaran yang sama,
dengan menurunkan identitas $\cos(x+T)=\cos(x)$ terhadap $x$, menghasilkan
$\sin(x+T)=\sin(x)$. Dengan demikian, setiap periode positif dari salah satu fungsi
memberikan, pada $x=0$, $\sin(T)=\sin(0)$ dan $\cos(T)=\cos(0)$.
Menurut rumus Euler, periode tersebut memenuhi
$e^{iT}=\cos(T)+i\sin(T)=\cos(0)+i\sin(0)=1$.
Hasil di atas menunjukkan bahwa tidak ada periode positif yang lebih kecil daripada $2\pi$.

Terakhir, kita ingin menunjukkan bahwa $e^{ix}$ merupakan pemetaan bijektif
dari himpunan $[0,2\pi)$ ke himpunan semua $z \in \C$ sedemikian sehingga
$\sabs{z} = 1$. Misalkan $e^{ix} = e^{iy}$ dan
$x > y$. Maka,
$0 < x-y < 2\pi$ karena $x,y \in [0,2\pi)$. Akan tetapi, $e^{i(x-y)} = 1$,
bertentangan dengan hasil di atas. Jadi, pemetaan tersebut injektif.
Untuk menunjukkan bahwa pemetaan $x \mapsto e^{ix}$ surjektif, ambil $(a,b) \in \R^2$ sedemikian sehingga $a^2+b^2 = 1$.
Pertama, andaikan $a,b \geq 0$. Menurut teorema nilai antara,
terdapat $x \in [0,\nicefrac{\pi}{2}]$ sedemikian sehingga
$\cos(x) = a$, dan akibatnya $b^2 = \bigl(\sin(x)\bigr)^2$. Karena
$b$ dan $\sin(x)$ nonnegatif, berlaku $b = \sin(x)$.
Karena $-\sin(x)$ merupakan turunan $\cos(x)$
dan $\cos(-x) = \cos(x)$,
kita memperoleh $\sin(x) < 0$ untuk $x \in [\nicefrac{-\pi}{2},0)$.
Dengan penalaran yang sama, kita memperoleh bahwa
jika $a > 0$ dan $b < 0$, kita dapat menemukan $x_0$ dalam $(\nicefrac{-\pi}{2},0)$
sedemikian sehingga $\cos(x_0) = a$ dan $\sin(x_0)=b$. Menurut periodisitas,
$x \coloneqq x_0+2\pi \in (\nicefrac{3\pi}{2},2\pi)$ juga memenuhi
$\cos(x) = a$ dan $\sin(x)=b$. Kasus $(a,b)=(0,-1)$ diperoleh dengan $x=\nicefrac{3\pi}{2}$.
Mengalikan dengan $-1$ sama dengan mengalikan dengan $e^{i\pi}$ atau
$e^{-i\pi}$. Jadi, kita selalu dapat mengasumsikan bahwa $a \geq 0$ (rincian diserahkan
sebagai latihan).
\end{proof}

\subsection{Lingkaran satuan dan koordinat polar}

Jika $\gamma \colon [a,b] \to \C$ merupakan kurva mulus sepotong-sepotong,
panjang busurnya diberikan oleh
\begin{equation*}
\int_a^b \babs{\gamma'(t)} \, dt .
\end{equation*}
Pemetaan $t \mapsto e^{it}$, untuk $t \in [0,2\pi)$, merupakan parametrisasi
bijektif lingkaran satuan. Untuk menghitung keliling, kita gunakan parametrisasi
satu putaran penuh pada interval tertutup $[0,2\pi]$. Karena
$\frac{d}{dt} \bigl( e^{it} \bigr) = ie^{it}$, keliling
lingkaran (panjang busurnya) adalah
\begin{equation*}
\int_0^{2\pi} \sabs{i e^{it}}  \,  dt
=
\int_0^{2\pi} 1  \,  dt  = 2\pi .
\end{equation*}

Lebih umum, pemetaan $t \mapsto e^{it}$ merupakan parametrisasi panjang busur bagi lingkaran.
Artinya, $t$, sebagai ukuran sudut dalam radian, mengukur panjang busur pada
lingkaran berjari-jari 1. Jadi, definisi $\sin$ dan $\cos$ yang diberikan
di atas sesuai dengan definisi geometris standar.

Setiap titik pada lingkaran satuan dapat dituliskan sebagai $e^{it}$ untuk suatu $t$.
Oleh karena itu, setiap bilangan kompleks $z \in \C$ dapat dituliskan dalam
\emph{\myindex{koordinat polar}} sebagai
\begin{equation*}
z = r e^{i\theta}
\end{equation*}
untuk suatu $r \geq 0$ dan $\theta \in \R$. Tentu saja, $\theta$
tidak tunggal karena $\theta$ dan $\theta+2\pi$ menghasilkan bilangan yang sama.
Hukum eksponen $e^{a+b} = e^a e^b$ menghasilkan rumus yang berguna untuk perpangkatan
dan hasil kali bilangan kompleks dalam koordinat polar. Rumus perpangkatan berlaku untuk
$n \in \N$, dan juga untuk $n \in \Z$ apabila $r > 0$:
\begin{equation*}
{(r e^{i\theta})}^n
= r^n e^{i n \theta} ,
\qquad
(r e^{i\theta})
(s e^{i\gamma})
=
rs e^{i(\theta+\gamma)} .
\end{equation*}

\subsection{Latihan}

\begin{exercise}
Turunkan deret pangkat untuk $\sin(z)$ dan $\cos(z)$ di titik asal.
\end{exercise}

\begin{exercise}
Dengan menggunakan deret pangkat, tunjukkan bahwa untuk $x$ real berlaku
$\frac{d}{dx} \bigl[ \sin(x)\bigr] = \cos(x)$ dan
$\frac{d}{dx} \bigl[ \cos(x)\bigr] = -\sin(x)$.
\end{exercise}

\begin{exercise}
Lengkapi bukti bahwa pemetaan $x \mapsto e^{ix}$ dari
$[0,2\pi)$ ke lingkaran satuan bersifat surjektif. Secara khusus, andaikan bahwa
kita telah memperoleh semua titik berbentuk $(a,b)$ dengan $a^2+b^2=1$ dan $a \geq 0$.
Dengan mengalikan dengan $e^{i\pi}$ atau $e^{-i\pi}$, tunjukkan bahwa kita memperoleh
semua titik, yakni juga titik-titik dengan $a < 0$.
\end{exercise}

\begin{exercise}
Buktikan bahwa fungsi eksponensial kompleks $z \mapsto e^z$ surjektif dari $\C$
ke $\C \setminus \{ 0 \}$, dan bahkan
bahwa untuk setiap $w$ tak nol, terdapat tak berhingga banyak $z \in \C$ sedemikian sehingga $e^z=w$.
\end{exercise}

\begin{exercise}
Buktikan bahwa untuk setiap $w \neq 0$ dan setiap $\epsilon > 0$,
terdapat $z \in \C$ dengan $\sabs{z} < \epsilon$ sedemikian sehingga $e^{1/z} = w$.
\end{exercise}

\begin{exercise}\label{exercise:cossinidentity}
Kita telah menunjukkan bahwa
${\bigl( \cos(x) \bigr)}^2 + {\bigl( \sin(x) \bigr)}^2 = 1$
untuk setiap $x \in \R$.
Buktikan bahwa
${\bigl( \cos(z) \bigr)}^2 + {\bigl( \sin(z) \bigr)}^2 = 1$
untuk setiap $z \in \C$.
\end{exercise}

\begin{exercise}
Buktikan identitas trigonometri
$\sin(z + w) = \sin(z) \cos(w) + \cos(z) \sin(w)$ dan
$\cos(z + w) = \cos(z) \cos(w) - \sin(z) \sin(w)$ untuk setiap $z,w \in \C$.
\end{exercise}

\begin{exercise}
Definisikan $\operatorname{sinc}(z) \coloneqq \frac{\sin(z)}{z}$ untuk $z \neq 0$ dan
$\operatorname{sinc}(0) \coloneqq 1$.
Tunjukkan bahwa fungsi sinc bersifat analitik dan hitung deret pangkatnya di titik asal.
\end{exercise}

\begin{exnote}
\pagebreak[2]
Definisikan \emph{\myindex{sinus hiperbolik}} dan
\emph{\myindex{kosinus hiperbolik}} melalui
\begin{equation*}
\sinh(z) \coloneqq \frac{e^z-e^{-z}}{2}, \qquad
\cosh(z) \coloneqq \frac{e^z+e^{-z}}{2}.
\end{equation*}
\end{exnote}

\begin{exercise}
Turunkan deret pangkat fungsi sinus hiperbolik dan fungsi kosinus hiperbolik di titik asal.
\end{exercise}

\begin{exercise}
Tunjukkan bahwa
\begin{enumerate}[a)]
\item
$\sinh(0) = 0$, $\cosh(0) = 1$.
\item
Untuk $x \in \R$,
$\frac{d}{dx} \bigl[ \sinh(x) \bigr] = \cosh(x)$ dan
$\frac{d}{dx} \bigl[ \cosh(x) \bigr] = \sinh(x)$.
\item
$\cosh(x) > 0$ untuk setiap $x \in \R$ dan
pemetaan $x \mapsto \sinh(x)$ naik tegas dan bijektif dari $\R$ ke $\R$.
\item
${\bigl(\cosh(x)\bigr)}^2 = 1 + {\bigl(\sinh(x)\bigr)}^2$ untuk setiap
$x$.
\end{enumerate}
\end{exercise}

\begin{exercise}
Definisikan $\tan(x) \coloneqq \frac{\sin(x)}{\cos(x)}$ untuk $x$ dengan $\cos(x) \neq 0$, seperti biasa.
\begin{enumerate}[a)]
\item
Tunjukkan bahwa untuk $x \in (\nicefrac{-\pi}{2},\nicefrac{\pi}{2})$,
baik $\sin$ maupun $\tan$ naik tegas, sehingga
fungsi invers $\arcsin$
dan $\arctan$ ada apabila kedua fungsi tersebut dibatasi pada interval tersebut.
\item
Tunjukkan bahwa $\arcsin$ diferensiabel pada $(-1,1)$ dan $\arctan$
diferensiabel pada $\R$, serta bahwa
$\frac{d}{dx} \arcsin(x) = \frac{1}{\sqrt{1-x^2}} \quad (-1<x<1)$ dan
$\frac{d}{dx} \arctan(x) = \frac{1}{1+x^2} \quad (x \in \R)$.
\item
Dengan menggunakan rumus jumlah geometri berhingga, tunjukkan bahwa deret pada
ruas kanan persamaan berikut konvergen untuk setiap $-1 \leq x \leq 1$
(termasuk titik-titik ujung), dan bahwa
\begin{equation*}
\arctan(x) = \int_0^x \frac{1}{1+t^2} dt
=
\sum_{k=0}^\infty \frac{{(-1)}^k}{2k+1} x^{2k+1}
\end{equation*}
Petunjuk: Integralkan jumlah berhingga tersebut, bukan deretnya.
\item
Gunakan hasil ini untuk menunjukkan bahwa
\begin{equation*}
1 - \frac{1}{3} + \frac{1}{5} - \cdots
=
\sum_{k=0}^\infty \frac{{(-1)}^k}{2k+1}
=
\frac{\pi}{4} .
\end{equation*}
\end{enumerate}
\end{exercise}

%%%%%%%%%%%%%%%%%%%%%%%%%%%%%%%%%%%%%%%%%%%%%%%%%%%%%%%%%%%%%%%%%%%%%%%%%%%%%%

% Deterministic reader cutoff after the admitted semantic boundary.
